\documentclass[
  doc,
  floatsintext,
  longtable,
  a4paper,
  nolmodern,
  notxfonts,
  notimes,
  12pt,
  colorlinks=true,linkcolor=blue,citecolor=blue,urlcolor=blue]{apa7}

\usepackage{amsmath}
\usepackage{amssymb}

\geometry{inner=1in, outer=1in}
\fancyhfoffset[LE,RO]{0cm}


\usepackage[bidi=default]{babel}
\babelprovide[main,import]{spanish}


% get rid of language-specific shorthands (see #6817):
\let\LanguageShortHands\languageshorthands
\def\languageshorthands#1{}

\RequirePackage{longtable}
\RequirePackage{threeparttablex}

\makeatletter
\renewcommand{\paragraph}{\@startsection{paragraph}{4}{\parindent}%
	{0\baselineskip \@plus 0.2ex \@minus 0.2ex}%
	{-.5em}%
	{\normalfont\normalsize\bfseries\typesectitle}}

\renewcommand{\subparagraph}[1]{\@startsection{subparagraph}{5}{0.5em}%
	{0\baselineskip \@plus 0.2ex \@minus 0.2ex}%
	{-\z@\relax}%
	{\normalfont\normalsize\bfseries\itshape\hspace{\parindent}{#1}\textit{\addperi}}{\relax}}
\makeatother




\usepackage{longtable, booktabs, multirow, multicol, colortbl, hhline, caption, array, float, xpatch}
\usepackage{subcaption}


\renewcommand\thesubfigure{\Alph{subfigure}}
\setcounter{topnumber}{2}
\setcounter{bottomnumber}{2}
\setcounter{totalnumber}{4}
\renewcommand{\topfraction}{0.85}
\renewcommand{\bottomfraction}{0.85}
\renewcommand{\textfraction}{0.15}
\renewcommand{\floatpagefraction}{0.7}

\usepackage{tcolorbox}
\tcbuselibrary{listings,theorems, breakable, skins}
\usepackage{fontawesome5}

\definecolor{quarto-callout-color}{HTML}{909090}
\definecolor{quarto-callout-note-color}{HTML}{0758E5}
\definecolor{quarto-callout-important-color}{HTML}{CC1914}
\definecolor{quarto-callout-warning-color}{HTML}{EB9113}
\definecolor{quarto-callout-tip-color}{HTML}{00A047}
\definecolor{quarto-callout-caution-color}{HTML}{FC5300}
\definecolor{quarto-callout-color-frame}{HTML}{ACACAC}
\definecolor{quarto-callout-note-color-frame}{HTML}{4582EC}
\definecolor{quarto-callout-important-color-frame}{HTML}{D9534F}
\definecolor{quarto-callout-warning-color-frame}{HTML}{F0AD4E}
\definecolor{quarto-callout-tip-color-frame}{HTML}{02B875}
\definecolor{quarto-callout-caution-color-frame}{HTML}{FD7E14}

%\newlength\Oldarrayrulewidth
%\newlength\Oldtabcolsep


\usepackage{hyperref}



\usepackage{color}
\usepackage{fancyvrb}
\newcommand{\VerbBar}{|}
\newcommand{\VERB}{\Verb[commandchars=\\\{\}]}
\DefineVerbatimEnvironment{Highlighting}{Verbatim}{commandchars=\\\{\}}
% Add ',fontsize=\small' for more characters per line
\usepackage{framed}
\definecolor{shadecolor}{RGB}{241,243,245}
\newenvironment{Shaded}{\begin{snugshade}}{\end{snugshade}}
\newcommand{\AlertTok}[1]{\textcolor[rgb]{0.68,0.00,0.00}{#1}}
\newcommand{\AnnotationTok}[1]{\textcolor[rgb]{0.37,0.37,0.37}{#1}}
\newcommand{\AttributeTok}[1]{\textcolor[rgb]{0.40,0.45,0.13}{#1}}
\newcommand{\BaseNTok}[1]{\textcolor[rgb]{0.68,0.00,0.00}{#1}}
\newcommand{\BuiltInTok}[1]{\textcolor[rgb]{0.00,0.23,0.31}{#1}}
\newcommand{\CharTok}[1]{\textcolor[rgb]{0.13,0.47,0.30}{#1}}
\newcommand{\CommentTok}[1]{\textcolor[rgb]{0.37,0.37,0.37}{#1}}
\newcommand{\CommentVarTok}[1]{\textcolor[rgb]{0.37,0.37,0.37}{\textit{#1}}}
\newcommand{\ConstantTok}[1]{\textcolor[rgb]{0.56,0.35,0.01}{#1}}
\newcommand{\ControlFlowTok}[1]{\textcolor[rgb]{0.00,0.23,0.31}{\textbf{#1}}}
\newcommand{\DataTypeTok}[1]{\textcolor[rgb]{0.68,0.00,0.00}{#1}}
\newcommand{\DecValTok}[1]{\textcolor[rgb]{0.68,0.00,0.00}{#1}}
\newcommand{\DocumentationTok}[1]{\textcolor[rgb]{0.37,0.37,0.37}{\textit{#1}}}
\newcommand{\ErrorTok}[1]{\textcolor[rgb]{0.68,0.00,0.00}{#1}}
\newcommand{\ExtensionTok}[1]{\textcolor[rgb]{0.00,0.23,0.31}{#1}}
\newcommand{\FloatTok}[1]{\textcolor[rgb]{0.68,0.00,0.00}{#1}}
\newcommand{\FunctionTok}[1]{\textcolor[rgb]{0.28,0.35,0.67}{#1}}
\newcommand{\ImportTok}[1]{\textcolor[rgb]{0.00,0.46,0.62}{#1}}
\newcommand{\InformationTok}[1]{\textcolor[rgb]{0.37,0.37,0.37}{#1}}
\newcommand{\KeywordTok}[1]{\textcolor[rgb]{0.00,0.23,0.31}{\textbf{#1}}}
\newcommand{\NormalTok}[1]{\textcolor[rgb]{0.00,0.23,0.31}{#1}}
\newcommand{\OperatorTok}[1]{\textcolor[rgb]{0.37,0.37,0.37}{#1}}
\newcommand{\OtherTok}[1]{\textcolor[rgb]{0.00,0.23,0.31}{#1}}
\newcommand{\PreprocessorTok}[1]{\textcolor[rgb]{0.68,0.00,0.00}{#1}}
\newcommand{\RegionMarkerTok}[1]{\textcolor[rgb]{0.00,0.23,0.31}{#1}}
\newcommand{\SpecialCharTok}[1]{\textcolor[rgb]{0.37,0.37,0.37}{#1}}
\newcommand{\SpecialStringTok}[1]{\textcolor[rgb]{0.13,0.47,0.30}{#1}}
\newcommand{\StringTok}[1]{\textcolor[rgb]{0.13,0.47,0.30}{#1}}
\newcommand{\VariableTok}[1]{\textcolor[rgb]{0.07,0.07,0.07}{#1}}
\newcommand{\VerbatimStringTok}[1]{\textcolor[rgb]{0.13,0.47,0.30}{#1}}
\newcommand{\WarningTok}[1]{\textcolor[rgb]{0.37,0.37,0.37}{\textit{#1}}}

\providecommand{\tightlist}{%
  \setlength{\itemsep}{0pt}\setlength{\parskip}{0pt}}
\usepackage{longtable,booktabs,array}
\usepackage{calc} % for calculating minipage widths
% Correct order of tables after \paragraph or \subparagraph
\usepackage{etoolbox}
\makeatletter
\patchcmd\longtable{\par}{\if@noskipsec\mbox{}\fi\par}{}{}
\makeatother
% Allow footnotes in longtable head/foot
\IfFileExists{footnotehyper.sty}{\usepackage{footnotehyper}}{\usepackage{footnote}}
\makesavenoteenv{longtable}

\usepackage{graphicx}
\makeatletter
\newsavebox\pandoc@box
\newcommand*\pandocbounded[1]{% scales image to fit in text height/width
  \sbox\pandoc@box{#1}%
  \Gscale@div\@tempa{\textheight}{\dimexpr\ht\pandoc@box+\dp\pandoc@box\relax}%
  \Gscale@div\@tempb{\linewidth}{\wd\pandoc@box}%
  \ifdim\@tempb\p@<\@tempa\p@\let\@tempa\@tempb\fi% select the smaller of both
  \ifdim\@tempa\p@<\p@\scalebox{\@tempa}{\usebox\pandoc@box}%
  \else\usebox{\pandoc@box}%
  \fi%
}
% Set default figure placement to htbp
\def\fps@figure{htbp}
\makeatother







\usepackage{newtx}

\defaultfontfeatures{Scale=MatchLowercase}
\defaultfontfeatures[\rmfamily]{Ligatures=TeX,Scale=1}





\title{Guía completa de comandos ADB Android}


\shorttitle{Guía completa de comandos ADB Android}


\usepackage{etoolbox}



\ccoppy{\textcopyright~2019}



\author{Edison Achalma}



\affiliation{
{Escuela Profesional de Economía, Universidad Nacional de San Cristóbal
de Huamanga}}




\leftheader{Achalma}

\date{2019-06-19}


\abstract{Este abstract será actualizado una vez que se complete el
contenido final del artículo. }

\keywords{keyword1, keyword2}

\authornote{\par{\addORCIDlink{Edison Achalma}{0000-0001-6996-3364}} 
\par{ }
\par{   El autor no tiene conflictos de interés que revelar.    Los
roles de autor se clasificaron utilizando la taxonomía de roles de
colaborador (CRediT; https://credit.niso.org/) de la siguiente
manera:  Edison Achalma:   conceptualización, metodología, análisis
formal, investigación, recursos, curación de
datos, redacción, visualización, supervisión, administración del
proyecto}
\par{La correspondencia relativa a este artículo debe dirigirse a Edison
Achalma, Escuela Profesional de Economía, Universidad Nacional de San
Cristóbal de Huamanga, Portal Independencia N°
57, Ayacucho, AYA 5001, Perú, Email: \href{mailto:elmer.achalma.09@unsch.edu.pe}{elmer.achalma.09@unsch.edu.pe}}
}

\makeatletter
\let\endoldlt\endlongtable
\def\endlongtable{
\hline
\endoldlt
}
\makeatother

\urlstyle{same}



\usepackage{hyperref}
\hypersetup{
  pdftitle={Guía completa de comandos ADB Android},
  pdfauthor={Edison Achalma},
  pdfsubject={Este abstract será actualizado una vez que se complete el contenido final del artículo.},
  pdfkeywords={keyword1, keyword2},
  unicode=true,
  colorlinks=true,
  linkcolor=blue,
  citecolor=blue,
  urlcolor=blue,
  pdfcreator={Quarto},
  pdflang={es}
}
\makeatletter
\@ifpackageloaded{caption}{}{\usepackage{caption}}
\AtBeginDocument{%
\ifdefined\contentsname
  \renewcommand*\contentsname{Tabla de contenidos}
\else
  \newcommand\contentsname{Tabla de contenidos}
\fi
\ifdefined\listfigurename
  \renewcommand*\listfigurename{Lista de Figuras}
\else
  \newcommand\listfigurename{Lista de Figuras}
\fi
\ifdefined\listtablename
  \renewcommand*\listtablename{Lista de Tablas}
\else
  \newcommand\listtablename{Lista de Tablas}
\fi
\ifdefined\figurename
  \renewcommand*\figurename{Figura}
\else
  \newcommand\figurename{Figura}
\fi
\ifdefined\tablename
  \renewcommand*\tablename{Tabla}
\else
  \newcommand\tablename{Tabla}
\fi
}
\@ifpackageloaded{float}{}{\usepackage{float}}
\floatstyle{ruled}
\@ifundefined{c@chapter}{\newfloat{codelisting}{h}{lop}}{\newfloat{codelisting}{h}{lop}[chapter]}
\floatname{codelisting}{Listado}
\newcommand*\listoflistings{\listof{codelisting}{Lista de Listados}}
\makeatother
\makeatletter
\makeatother
\makeatletter
\@ifpackageloaded{caption}{}{\usepackage{caption}}
\@ifpackageloaded{subcaption}{}{\usepackage{subcaption}}
\makeatother
\makeatletter
\@ifpackageloaded{fontawesome5}{}{\usepackage{fontawesome5}}
\makeatother

% From https://tex.stackexchange.com/a/645996/211326
%%% apa7 doesn't want to add appendix section titles in the toc
%%% let's make it do it
\makeatletter
\xpatchcmd{\appendix}
  {\par}
  {\addcontentsline{toc}{section}{\@currentlabelname}\par}
  {}{}
\makeatother

%% Disable longtable counter
%% https://tex.stackexchange.com/a/248395/211326

\usepackage{etoolbox}

\makeatletter
\patchcmd{\LT@caption}
  {\bgroup}
  {\bgroup\global\LTpatch@captiontrue}
  {}{}
\patchcmd{\longtable}
  {\par}
  {\par\global\LTpatch@captionfalse}
  {}{}
\apptocmd{\endlongtable}
  {\ifLTpatch@caption\else\addtocounter{table}{-1}\fi}
  {}{}
\newif\ifLTpatch@caption
\makeatother

\begin{document}

\maketitle


\hypertarget{toc}{}
\tableofcontents
\newpage
\section[Introduction]{Guía completa de comandos ADB Android}

\setcounter{secnumdepth}{3}

\setlength\LTleft{0pt}




\section{ADB}\label{adb}

\subsection{¿Qué es Android Debug
Bridge?}\label{quuxe9-es-android-debug-bridge}

\textbf{Android Debug Bridge (ADB)} es una herramienta de línea de
comandos versátil que permite comunicarte con dispositivos Android. Es
parte esencial del SDK de Android y proporciona acceso completo al
sistema operativo del dispositivo.

\subsection{Componentes de ADB}\label{componentes-de-adb}

ADB es un programa cliente-servidor con tres componentes:

\begin{enumerate}
\def\labelenumi{\arabic{enumi}.}
\item
  \textbf{Cliente}

  \begin{itemize}
  \tightlist
  \item
    Se ejecuta en tu computadora (Arch Linux)
  \item
    Envía comandos al servidor
  \item
    Se invoca desde la terminal con \texttt{adb}
  \end{itemize}
\item
  \textbf{Daemon (adbd)}

  \begin{itemize}
  \tightlist
  \item
    Se ejecuta en el dispositivo Android
  \item
    Ejecuta los comandos recibidos
  \item
    Proceso en segundo plano
  \end{itemize}
\item
  \textbf{Servidor}

  \begin{itemize}
  \tightlist
  \item
    Se ejecuta en tu computadora
  \item
    Gestiona la comunicación entre cliente y daemon
  \item
    Puerto TCP 5037
  \item
    Proceso en segundo plano
  \end{itemize}
\end{enumerate}

\subsection{Funcionamiento}\label{funcionamiento}

\begin{verbatim}
┌─────────────────────────────────────────────┐
│          ARCH LINUX (Host)                  │
│                                             │
│  ┌──────────┐         ┌──────────┐         │
│  │ Cliente  │────────▶│ Servidor │         │
│  │   adb    │         │ (puerto  │         │
│  │          │         │  5037)   │         │
│  └──────────┘         └─────┬────┘         │
│                             │              │
└─────────────────────────────┼──────────────┘
                              │
                              │ USB / WiFi
                              │
┌─────────────────────────────▼──────────────┐
│       DISPOSITIVO ANDROID                  │
│                                            │
│         ┌──────────────┐                   │
│         │   Daemon     │                   │
│         │   (adbd)     │                   │
│         └──────────────┘                   │
│                                            │
└────────────────────────────────────────────┘
\end{verbatim}

\subsection{Ventajas de ADB}\label{ventajas-de-adb}

\textbf{Para desarrollo:}

\begin{itemize}
\tightlist
\item
  Instalar y depurar aplicaciones
\item
  Acceder al shell Unix del dispositivo
\item
  Copiar archivos bidireccional
\item
  Ejecutar comandos del sistema
\end{itemize}

\textbf{Para testing:}

\begin{itemize}
\tightlist
\item
  Automatizar pruebas
\item
  Capturar logs y errores
\item
  Simular eventos del usuario
\item
  Modificar configuraciones
\end{itemize}

\textbf{Para gestión:}

\begin{itemize}
\tightlist
\item
  Instalar APKs sin Google Play
\item
  Hacer backups de datos
\item
  Capturar pantalla y video
\item
  Analizar rendimiento
\end{itemize}

\begin{center}\rule{0.5\linewidth}{0.5pt}\end{center}

\section{Arquitectura de ADB}\label{arquitectura-de-adb}

\subsection{Proceso de Conexión}\label{proceso-de-conexiuxf3n}

\textbf{1. Inicio del Cliente}

\begin{Shaded}
\begin{Highlighting}[]
\ExtensionTok{adb}\NormalTok{ devices}
\end{Highlighting}
\end{Shaded}

\textbf{2. Verificación del Servidor}

\begin{itemize}
\tightlist
\item
  Si no hay servidor corriendo → lo inicia
\item
  El servidor se vincula al puerto 5037
\item
  Escucha comandos de clientes
\end{itemize}

\textbf{3. Detección de Dispositivos}

\begin{itemize}
\tightlist
\item
  Escanea puertos USB
\item
  Escanea puertos de red (5555-5585)
\item
  Establece conexión con daemon (adbd)
\end{itemize}

\textbf{4. Comunicación}

\begin{itemize}
\tightlist
\item
  Cliente envía comando al servidor
\item
  Servidor envía comando al daemon
\item
  Daemon ejecuta y responde
\item
  Respuesta llega al cliente
\end{itemize}

\subsection{Puertos y Conexiones}\label{puertos-y-conexiones}

\textbf{Emuladores:}

\begin{verbatim}
Emulador 1: consola 5554, adb 5555
Emulador 2: consola 5556, adb 5557
Emulador 3: consola 5558, adb 5559
...hasta 16 emuladores
\end{verbatim}

\textbf{Dispositivos Físicos:}

\begin{itemize}
\tightlist
\item
  USB: Detección automática por vendor ID
\item
  WiFi: Puerto 5555 (configurable)
\end{itemize}

\section{Instalación en Arch Linux}\label{instalaciuxf3n-en-arch-linux}

\subsection{Método 1: Repositorios Oficiales
(Recomendado)}\label{muxe9todo-1-repositorios-oficiales-recomendado}

\begin{Shaded}
\begin{Highlighting}[]
\CommentTok{\# ADB está en el paquete android{-}tools}
\FunctionTok{sudo}\NormalTok{ pacman }\AttributeTok{{-}S}\NormalTok{ android{-}tools}
\end{Highlighting}
\end{Shaded}

\textbf{Verificar instalación:}

\begin{Shaded}
\begin{Highlighting}[]
\ExtensionTok{adb}\NormalTok{ version}
\CommentTok{\# Salida esperada:}
\CommentTok{\# Android Debug Bridge version 1.0.41}
\CommentTok{\# Version 35.0.1{-}11580240}
\end{Highlighting}
\end{Shaded}

\subsection{Método 2: Android SDK
Completo}\label{muxe9todo-2-android-sdk-completo}

\textbf{Si necesitas todas las herramientas de Android:}

\begin{Shaded}
\begin{Highlighting}[]
\CommentTok{\# 1. Instalar dependencias}
\FunctionTok{sudo}\NormalTok{ pacman }\AttributeTok{{-}S}\NormalTok{ jdk17{-}openjdk}

\CommentTok{\# 2. Instalar Android SDK desde AUR}
\ExtensionTok{yay} \AttributeTok{{-}S}\NormalTok{ android{-}sdk android{-}sdk{-}platform{-}tools}

\CommentTok{\# 3. Configurar variables de entorno}
\CommentTok{\# Agregar a \textasciitilde{}/.config/fish/config.fish}
\BuiltInTok{set} \AttributeTok{{-}gx}\NormalTok{ ANDROID\_HOME /opt/android{-}sdk}
\BuiltInTok{set} \AttributeTok{{-}gx}\NormalTok{ PATH }\VariableTok{$PATH} \VariableTok{$ANDROID\_HOME}\NormalTok{/platform{-}tools }\VariableTok{$ANDROID\_HOME}\NormalTok{/tools}
\end{Highlighting}
\end{Shaded}

\textbf{Recargar configuración:}

\begin{Shaded}
\begin{Highlighting}[]
\NormalTok{source \textasciitilde{}/.config/fish/config.fish}
\end{Highlighting}
\end{Shaded}

\subsection{Método 3: Descarga
Manual}\label{muxe9todo-3-descarga-manual}

\begin{Shaded}
\begin{Highlighting}[]
\CommentTok{\# 1. Crear directorio}
\FunctionTok{mkdir} \AttributeTok{{-}p}\NormalTok{ \textasciitilde{}/android/platform{-}tools}
\BuiltInTok{cd}\NormalTok{ \textasciitilde{}/android/platform{-}tools}

\CommentTok{\# 2. Descargar herramientas}
\ExtensionTok{curl} \AttributeTok{{-}LO}\NormalTok{ https://dl.google.com/android/repository/platform{-}tools{-}latest{-}linux.zip}

\CommentTok{\# 3. Extraer}
\FunctionTok{unzip}\NormalTok{ platform{-}tools{-}latest{-}linux.zip}

\CommentTok{\# 4. Agregar al PATH (Fish)}
\BuiltInTok{set} \AttributeTok{{-}gx}\NormalTok{ PATH }\VariableTok{$PATH}\NormalTok{ \textasciitilde{}/android/platform{-}tools}

\CommentTok{\# Para hacerlo permanente:}
\BuiltInTok{echo} \StringTok{\textquotesingle{}set {-}gx PATH $PATH \textasciitilde{}/android/platform{-}tools\textquotesingle{}} \OperatorTok{\textgreater{}\textgreater{}}\NormalTok{ \textasciitilde{}/.config/fish/config.fish}
\end{Highlighting}
\end{Shaded}

\section{Configuración Inicial}\label{configuraciuxf3n-inicial}

\subsection{1. Verificar Instalación}\label{verificar-instalaciuxf3n}

\begin{Shaded}
\begin{Highlighting}[]
\CommentTok{\# Versión de ADB}
\ExtensionTok{adb}\NormalTok{ version}

\CommentTok{\# Ayuda general}
\ExtensionTok{adb} \AttributeTok{{-}{-}help}

\CommentTok{\# Lista de comandos}
\ExtensionTok{adb}\NormalTok{ help}
\end{Highlighting}
\end{Shaded}

\subsection{2. Configurar Permisos de
Usuario}\label{configurar-permisos-de-usuario}

\textbf{En Arch Linux, necesitas estar en el grupo \texttt{adbusers}:}

\begin{Shaded}
\begin{Highlighting}[]
\CommentTok{\# Verificar si el grupo existe}
\FunctionTok{getent}\NormalTok{ group adbusers}

\CommentTok{\# Si no existe, crearlo}
\FunctionTok{sudo}\NormalTok{ groupadd adbusers}

\CommentTok{\# Agregar tu usuario al grupo}
\FunctionTok{sudo}\NormalTok{ usermod }\AttributeTok{{-}aG}\NormalTok{ adbusers }\VariableTok{$USER}

\CommentTok{\# Verificar}
\FunctionTok{groups} \VariableTok{$USER} \KeywordTok{|} \FunctionTok{grep}\NormalTok{ adbusers}
\end{Highlighting}
\end{Shaded}

\textbf{Aplicar cambios (una de estas opciones):}

\begin{Shaded}
\begin{Highlighting}[]
\CommentTok{\# Opción 1: Recargar grupos (temporal)}
\ExtensionTok{newgrp}\NormalTok{ adbusers}

\CommentTok{\# Opción 2: Cerrar sesión y volver a iniciar}

\CommentTok{\# Opción 3: Reiniciar sistema}
\end{Highlighting}
\end{Shaded}

\subsection{3. Configurar Reglas udev}\label{configurar-reglas-udev}

\textbf{Las reglas udev permiten que ADB detecte dispositivos sin sudo.}

\begin{Shaded}
\begin{Highlighting}[]
\CommentTok{\# Crear archivo de reglas}
\FunctionTok{sudo}\NormalTok{ nano /etc/udev/rules.d/51{-}android.rules}
\end{Highlighting}
\end{Shaded}

\textbf{Agregar reglas para fabricantes comunes:}

\begin{verbatim}
# Google
SUBSYSTEM=="usb", ATTR{idVendor}=="18d1", MODE="0666", GROUP="adbusers"

# Samsung
SUBSYSTEM=="usb", ATTR{idVendor}=="04e8", MODE="0666", GROUP="adbusers"

# Xiaomi
SUBSYSTEM=="usb", ATTR{idVendor}=="2717", MODE="0666", GROUP="adbusers"

# Huawei
SUBSYSTEM=="usb", ATTR{idVendor}=="12d1", MODE="0666", GROUP="adbusers"

# OnePlus
SUBSYSTEM=="usb", ATTR{idVendor}=="2a70", MODE="0666", GROUP="adbusers"

# Motorola
SUBSYSTEM=="usb", ATTR{idVendor}=="22b8", MODE="0666", GROUP="adbusers"

# LG
SUBSYSTEM=="usb", ATTR{idVendor}=="1004", MODE="0666", GROUP="adbusers"

# HTC
SUBSYSTEM=="usb", ATTR{idVendor}=="0bb4", MODE="0666", GROUP="adbusers"

# Sony
SUBSYSTEM=="usb", ATTR{idVendor}=="0fce", MODE="0666", GROUP="adbusers"

# ASUS
SUBSYSTEM=="usb", ATTR{idVendor}=="0b05", MODE="0666", GROUP="adbusers"

# Oppo
SUBSYSTEM=="usb", ATTR{idVendor}=="22d9", MODE="0666", GROUP="adbusers"

# Realme
SUBSYSTEM=="usb", ATTR{idVendor}=="22d9", MODE="0666", GROUP="adbusers"

# Vivo
SUBSYSTEM=="usb", ATTR{idVendor}=="2d95", MODE="0666", GROUP="adbusers"
\end{verbatim}

\textbf{Recargar reglas:}

\begin{Shaded}
\begin{Highlighting}[]
\FunctionTok{sudo}\NormalTok{ udevadm control }\AttributeTok{{-}{-}reload{-}rules}
\FunctionTok{sudo}\NormalTok{ udevadm trigger}
\end{Highlighting}
\end{Shaded}

\subsection{4. Habilitar Depuración USB en
Android}\label{habilitar-depuraciuxf3n-usb-en-android}

\textbf{En el dispositivo Android:}

\begin{enumerate}
\def\labelenumi{\arabic{enumi}.}
\item
  \textbf{Activar Opciones de Desarrollador:}

  \begin{itemize}
  \tightlist
  \item
    Ir a: \texttt{Ajustes} → \texttt{Acerca\ del\ teléfono}
  \item
    Tocar 7 veces en \texttt{Número\ de\ compilación}
  \end{itemize}
\item
  \textbf{Habilitar Depuración USB:}

  \begin{itemize}
  \tightlist
  \item
    Ir a: \texttt{Ajustes} → \texttt{Opciones\ de\ desarrollador}
  \item
    Activar: \texttt{Depuración\ USB}
  \end{itemize}
\item
  \textbf{Dispositivos Xiaomi (adicional):}

  \begin{itemize}
  \tightlist
  \item
    Activar: \texttt{Depuración\ USB\ (Configuración\ de\ seguridad)}
  \item
    Reiniciar dispositivo
  \end{itemize}
\end{enumerate}

\subsection{5. Primera Conexión}\label{primera-conexiuxf3n}

\begin{Shaded}
\begin{Highlighting}[]
\CommentTok{\# Conectar dispositivo por USB}

\CommentTok{\# Iniciar servidor ADB}
\ExtensionTok{adb}\NormalTok{ start{-}server}

\CommentTok{\# Listar dispositivos}
\ExtensionTok{adb}\NormalTok{ devices}
\end{Highlighting}
\end{Shaded}

\textbf{En el dispositivo Android aparecerá un diálogo:}

\begin{itemize}
\tightlist
\item
  ``¿Permitir depuración USB?''
\item
  Marcar: ``Permitir siempre desde esta computadora''
\item
  Aceptar
\end{itemize}

\textbf{Verificar conexión:}

\begin{Shaded}
\begin{Highlighting}[]
\ExtensionTok{adb}\NormalTok{ devices}
\CommentTok{\# Salida esperada:}
\CommentTok{\# List of devices attached}
\CommentTok{\# ABC123456789    device}
\end{Highlighting}
\end{Shaded}

\section{Conexión de Dispositivos}\label{conexiuxf3n-de-dispositivos}

\subsection{Conexión USB}\label{conexiuxf3n-usb}

\textbf{Básica (un solo dispositivo):}

\begin{Shaded}
\begin{Highlighting}[]
\CommentTok{\# Listar dispositivos}
\ExtensionTok{adb}\NormalTok{ devices}

\CommentTok{\# Salida:}
\CommentTok{\# List of devices attached}
\CommentTok{\# ABC123456789    device}
\end{Highlighting}
\end{Shaded}

\textbf{Múltiples dispositivos:}

\begin{Shaded}
\begin{Highlighting}[]
\CommentTok{\# Listar con detalles}
\ExtensionTok{adb}\NormalTok{ devices }\AttributeTok{{-}l}

\CommentTok{\# Salida:}
\CommentTok{\# List of devices attached}
\CommentTok{\# ABC123    device usb:1{-}1 product:phone model:Pixel\_6 device:pixel6}
\CommentTok{\# DEF456    device usb:1{-}2 product:tablet model:Tab\_S8 device:tabs8}
\end{Highlighting}
\end{Shaded}

\textbf{Especificar dispositivo:}

\begin{Shaded}
\begin{Highlighting}[]
\CommentTok{\# Por número de serie}
\ExtensionTok{adb} \AttributeTok{{-}s}\NormalTok{ ABC123456789 shell}

\CommentTok{\# Solo USB (si hay uno solo)}
\ExtensionTok{adb} \AttributeTok{{-}d}\NormalTok{ shell}

\CommentTok{\# Solo emulador (si hay uno solo)}
\ExtensionTok{adb} \AttributeTok{{-}e}\NormalTok{ shell}
\end{Highlighting}
\end{Shaded}

\subsection{Conexión WiFi (Android
11+)}\label{conexiuxf3n-wifi-android-11}

\subsubsection{Método 1: Vinculación con
Código}\label{muxe9todo-1-vinculaciuxf3n-con-cuxf3digo}

\textbf{En el dispositivo:}

\begin{enumerate}
\def\labelenumi{\arabic{enumi}.}
\tightlist
\item
  Ir a: \texttt{Ajustes} → \texttt{Opciones\ de\ desarrollador}
\item
  Activar: \texttt{Depuración\ inalámbrica}
\item
  Tocar: \texttt{Vincular\ dispositivo\ con\ código\ de\ vinculación}
\item
  Aparecerá: IP, Puerto y Código de 6 dígitos
\end{enumerate}

\textbf{En Arch Linux:}

\begin{Shaded}
\begin{Highlighting}[]
\CommentTok{\# Vincular (usar IP y puerto del dispositivo)}
\ExtensionTok{adb}\NormalTok{ pair 192.168.1.100:37859}

\CommentTok{\# Ingresar código de 6 dígitos cuando se solicite}
\CommentTok{\# Enter pairing code: 123456}

\CommentTok{\# Conectar (usar IP mostrada en "Depuración inalámbrica")}
\ExtensionTok{adb}\NormalTok{ connect 192.168.1.100:5555}

\CommentTok{\# Verificar}
\ExtensionTok{adb}\NormalTok{ devices}
\end{Highlighting}
\end{Shaded}

\subsubsection{Método 2: Desde USB a WiFi (Android 10 y
anteriores)}\label{muxe9todo-2-desde-usb-a-wifi-android-10-y-anteriores}

\begin{Shaded}
\begin{Highlighting}[]
\CommentTok{\# 1. Conectar dispositivo por USB}
\ExtensionTok{adb}\NormalTok{ devices}

\CommentTok{\# 2. Obtener IP del dispositivo}
\ExtensionTok{adb}\NormalTok{ shell ip addr show wlan0 }\KeywordTok{|} \FunctionTok{grep}\NormalTok{ inet}
\CommentTok{\# O desde el dispositivo: Ajustes → Acerca del teléfono → Estado}

\CommentTok{\# 3. Habilitar TCP/IP en puerto 5555}
\ExtensionTok{adb}\NormalTok{ tcpip 5555}

\CommentTok{\# 4. Desconectar cable USB}

\CommentTok{\# 5. Conectar por WiFi}
\ExtensionTok{adb}\NormalTok{ connect 192.168.1.100:5555}

\CommentTok{\# 6. Verificar}
\ExtensionTok{adb}\NormalTok{ devices}
\end{Highlighting}
\end{Shaded}

\textbf{Volver a USB:}

\begin{Shaded}
\begin{Highlighting}[]
\ExtensionTok{adb}\NormalTok{ usb}
\end{Highlighting}
\end{Shaded}

\subsection{Reconexión Automática}\label{reconexiuxf3n-automuxe1tica}

\textbf{Si se pierde conexión WiFi:}

\begin{Shaded}
\begin{Highlighting}[]
\CommentTok{\# Reintentar conexión}
\ExtensionTok{adb}\NormalTok{ connect 192.168.1.100:5555}

\CommentTok{\# Si no funciona, reiniciar servidor}
\ExtensionTok{adb}\NormalTok{ kill{-}server}
\ExtensionTok{adb}\NormalTok{ start{-}server}
\ExtensionTok{adb}\NormalTok{ connect 192.168.1.100:5555}
\end{Highlighting}
\end{Shaded}

\subsection{Variable de Entorno
ANDROID\_SERIAL}\label{variable-de-entorno-android_serial}

\begin{Shaded}
\begin{Highlighting}[]
\CommentTok{\# Fish}
\BuiltInTok{set} \AttributeTok{{-}gx}\NormalTok{ ANDROID\_SERIAL ABC123456789}
\ExtensionTok{adb}\NormalTok{ shell  }\CommentTok{\# Se conecta automáticamente a ABC123456789}

\CommentTok{\# Bash/Zsh}
\BuiltInTok{export} \VariableTok{ANDROID\_SERIAL}\OperatorTok{=}\NormalTok{ABC123456789}
\ExtensionTok{adb}\NormalTok{ shell}

\CommentTok{\# Hacer permanente en Fish}
\BuiltInTok{echo} \StringTok{\textquotesingle{}set {-}gx ANDROID\_SERIAL ABC123456789\textquotesingle{}} \OperatorTok{\textgreater{}\textgreater{}}\NormalTok{ \textasciitilde{}/.config/fish/config.fish}
\end{Highlighting}
\end{Shaded}

\begin{center}\rule{0.5\linewidth}{0.5pt}\end{center}

\section{Comandos Básicos}\label{comandos-buxe1sicos}

\subsection{Gestión del Servidor ADB}\label{gestiuxf3n-del-servidor-adb}

\begin{Shaded}
\begin{Highlighting}[]
\CommentTok{\# Iniciar servidor}
\ExtensionTok{adb}\NormalTok{ start{-}server}

\CommentTok{\# Detener servidor}
\ExtensionTok{adb}\NormalTok{ kill{-}server}

\CommentTok{\# Reiniciar servidor}
\ExtensionTok{adb}\NormalTok{ kill{-}server }\KeywordTok{\&\&} \ExtensionTok{adb}\NormalTok{ start{-}server}

\CommentTok{\# Verificar versión}
\ExtensionTok{adb}\NormalTok{ version}

\CommentTok{\# Ver puerto del servidor}
\CommentTok{\# (siempre es 5037 por defecto)}
\end{Highlighting}
\end{Shaded}

\subsection{Información del
Dispositivo}\label{informaciuxf3n-del-dispositivo}

\begin{Shaded}
\begin{Highlighting}[]
\CommentTok{\# Listar dispositivos conectados}
\ExtensionTok{adb}\NormalTok{ devices}
\ExtensionTok{adb}\NormalTok{ devices }\AttributeTok{{-}l}  \CommentTok{\# con detalles}

\CommentTok{\# Shell del dispositivo}
\ExtensionTok{adb}\NormalTok{ shell}

\CommentTok{\# Ejecutar comando único}
\ExtensionTok{adb}\NormalTok{ shell ls /sdcard}

\CommentTok{\# Obtener información del sistema}
\ExtensionTok{adb}\NormalTok{ shell getprop}

\CommentTok{\# Versión de Android}
\ExtensionTok{adb}\NormalTok{ shell getprop ro.build.version.release}

\CommentTok{\# Modelo del dispositivo}
\ExtensionTok{adb}\NormalTok{ shell getprop ro.product.model}

\CommentTok{\# Fabricante}
\ExtensionTok{adb}\NormalTok{ shell getprop ro.product.manufacturer}

\CommentTok{\# Número de serie}
\ExtensionTok{adb}\NormalTok{ get{-}serialno}

\CommentTok{\# Estado del dispositivo}
\ExtensionTok{adb}\NormalTok{ get{-}state}
\CommentTok{\# Posibles valores: device, offline, bootloader, recovery}
\end{Highlighting}
\end{Shaded}

\subsection{Reiniciar Dispositivo}\label{reiniciar-dispositivo}

\begin{Shaded}
\begin{Highlighting}[]
\CommentTok{\# Reinicio normal}
\ExtensionTok{adb}\NormalTok{ reboot}

\CommentTok{\# Reiniciar en recovery}
\ExtensionTok{adb}\NormalTok{ reboot recovery}

\CommentTok{\# Reiniciar en bootloader}
\ExtensionTok{adb}\NormalTok{ reboot bootloader}

\CommentTok{\# Reiniciar en modo fastboot}
\ExtensionTok{adb}\NormalTok{ reboot{-}bootloader}
\end{Highlighting}
\end{Shaded}

\subsection{Logs y Debugging}\label{logs-y-debugging}

\begin{Shaded}
\begin{Highlighting}[]
\CommentTok{\# Ver logs en tiempo real}
\ExtensionTok{adb}\NormalTok{ logcat}

\CommentTok{\# Filtrar por tag}
\ExtensionTok{adb}\NormalTok{ logcat ActivityManager:I }\PreprocessorTok{*}\NormalTok{:S}

\CommentTok{\# Guardar logs en archivo}
\ExtensionTok{adb}\NormalTok{ logcat }\OperatorTok{\textgreater{}}\NormalTok{ logs.txt}

\CommentTok{\# Limpiar buffer de logs}
\ExtensionTok{adb}\NormalTok{ logcat }\AttributeTok{{-}c}

\CommentTok{\# Ver logs de kernel}
\ExtensionTok{adb}\NormalTok{ shell dmesg}

\CommentTok{\# Obtener bug report}
\ExtensionTok{adb}\NormalTok{ bugreport bugreport.zip}
\end{Highlighting}
\end{Shaded}

\begin{center}\rule{0.5\linewidth}{0.5pt}\end{center}

\section{Gestión de Aplicaciones}\label{gestiuxf3n-de-aplicaciones}

\subsection{Instalar APK}\label{instalar-apk}

\begin{Shaded}
\begin{Highlighting}[]
\CommentTok{\# Instalación básica}
\ExtensionTok{adb}\NormalTok{ install app.apk}

\CommentTok{\# Reinstalar manteniendo datos}
\ExtensionTok{adb}\NormalTok{ install }\AttributeTok{{-}r}\NormalTok{ app.apk}

\CommentTok{\# Instalar en almacenamiento externo}
\ExtensionTok{adb}\NormalTok{ install }\AttributeTok{{-}s}\NormalTok{ app.apk}

\CommentTok{\# Permitir downgrade}
\ExtensionTok{adb}\NormalTok{ install }\AttributeTok{{-}d}\NormalTok{ app.apk}

\CommentTok{\# Instalar APK de prueba}
\ExtensionTok{adb}\NormalTok{ install }\AttributeTok{{-}t}\NormalTok{ app.apk}

\CommentTok{\# Otorgar todos los permisos}
\ExtensionTok{adb}\NormalTok{ install }\AttributeTok{{-}g}\NormalTok{ app.apk}

\CommentTok{\# Instalar en dispositivo específico}
\ExtensionTok{adb} \AttributeTok{{-}s}\NormalTok{ ABC123 install app.apk}

\CommentTok{\# Instalar múltiples APKs (split APKs)}
\ExtensionTok{adb}\NormalTok{ install{-}multiple base.apk config.apk}

\CommentTok{\# Instalación incremental (Android 11+)}
\ExtensionTok{adb}\NormalTok{ install }\AttributeTok{{-}{-}incremental}\NormalTok{ app.apk}
\end{Highlighting}
\end{Shaded}

\subsection{Desinstalar Aplicación}\label{desinstalar-aplicaciuxf3n}

\begin{Shaded}
\begin{Highlighting}[]
\CommentTok{\# Desinstalar}
\ExtensionTok{adb}\NormalTok{ uninstall com.example.app}

\CommentTok{\# Desinstalar pero mantener datos}
\ExtensionTok{adb}\NormalTok{ uninstall }\AttributeTok{{-}k}\NormalTok{ com.example.app}

\CommentTok{\# Desinstalar del dispositivo específico}
\ExtensionTok{adb} \AttributeTok{{-}s}\NormalTok{ ABC123 uninstall com.example.app}
\end{Highlighting}
\end{Shaded}

\subsection{Listar Aplicaciones}\label{listar-aplicaciones}

\begin{Shaded}
\begin{Highlighting}[]
\CommentTok{\# Todas las aplicaciones}
\ExtensionTok{adb}\NormalTok{ shell pm list packages}

\CommentTok{\# Solo aplicaciones de terceros}
\ExtensionTok{adb}\NormalTok{ shell pm list packages }\AttributeTok{{-}3}

\CommentTok{\# Solo aplicaciones del sistema}
\ExtensionTok{adb}\NormalTok{ shell pm list packages }\AttributeTok{{-}s}

\CommentTok{\# Aplicaciones habilitadas}
\ExtensionTok{adb}\NormalTok{ shell pm list packages }\AttributeTok{{-}e}

\CommentTok{\# Aplicaciones deshabilitadas}
\ExtensionTok{adb}\NormalTok{ shell pm list packages }\AttributeTok{{-}d}

\CommentTok{\# Buscar por nombre}
\ExtensionTok{adb}\NormalTok{ shell pm list packages }\KeywordTok{|} \FunctionTok{grep}\NormalTok{ chrome}

\CommentTok{\# Con ruta del APK}
\ExtensionTok{adb}\NormalTok{ shell pm list packages }\AttributeTok{{-}f}
\end{Highlighting}
\end{Shaded}

\subsection{Información de
Aplicación}\label{informaciuxf3n-de-aplicaciuxf3n}

\begin{Shaded}
\begin{Highlighting}[]
\CommentTok{\# Ruta del APK}
\ExtensionTok{adb}\NormalTok{ shell pm path com.android.chrome}

\CommentTok{\# Información completa}
\ExtensionTok{adb}\NormalTok{ shell dumpsys package com.android.chrome}

\CommentTok{\# Permisos}
\ExtensionTok{adb}\NormalTok{ shell dumpsys package com.android.chrome }\KeywordTok{|} \FunctionTok{grep}\NormalTok{ permission}

\CommentTok{\# Versión}
\ExtensionTok{adb}\NormalTok{ shell dumpsys package com.android.chrome }\KeywordTok{|} \FunctionTok{grep}\NormalTok{ versionName}
\end{Highlighting}
\end{Shaded}

\subsection{Iniciar Aplicación}\label{iniciar-aplicaciuxf3n}

\begin{Shaded}
\begin{Highlighting}[]
\CommentTok{\# Iniciar actividad principal}
\ExtensionTok{adb}\NormalTok{ shell monkey }\AttributeTok{{-}p}\NormalTok{ com.android.chrome }\AttributeTok{{-}c}\NormalTok{ android.intent.category.LAUNCHER 1}

\CommentTok{\# Iniciar actividad específica}
\ExtensionTok{adb}\NormalTok{ shell am start }\AttributeTok{{-}n}\NormalTok{ com.android.chrome/com.google.android.apps.chrome.Main}

\CommentTok{\# Iniciar con intent}
\ExtensionTok{adb}\NormalTok{ shell am start }\AttributeTok{{-}a}\NormalTok{ android.intent.action.VIEW }\AttributeTok{{-}d}\NormalTok{ https://google.com}
\end{Highlighting}
\end{Shaded}

\subsection{Detener Aplicación}\label{detener-aplicaciuxf3n}

\begin{Shaded}
\begin{Highlighting}[]
\CommentTok{\# Forzar detención}
\ExtensionTok{adb}\NormalTok{ shell am force{-}stop com.android.chrome}

\CommentTok{\# Matar proceso}
\ExtensionTok{adb}\NormalTok{ shell am kill com.android.chrome}
\end{Highlighting}
\end{Shaded}

\subsection{Limpiar Datos de
Aplicación}\label{limpiar-datos-de-aplicaciuxf3n}

\begin{Shaded}
\begin{Highlighting}[]
\CommentTok{\# Borrar todos los datos y caché}
\ExtensionTok{adb}\NormalTok{ shell pm clear com.android.chrome}
\end{Highlighting}
\end{Shaded}

\subsection{Habilitar/Deshabilitar
Aplicación}\label{habilitardeshabilitar-aplicaciuxf3n}

\begin{Shaded}
\begin{Highlighting}[]
\CommentTok{\# Deshabilitar}
\ExtensionTok{adb}\NormalTok{ shell pm disable{-}user com.example.app}

\CommentTok{\# Habilitar}
\ExtensionTok{adb}\NormalTok{ shell pm enable com.example.app}
\end{Highlighting}
\end{Shaded}

\section{Transferencia de Archivos}\label{transferencia-de-archivos}

\subsection{Copiar Archivos al Dispositivo
(Push)}\label{copiar-archivos-al-dispositivo-push}

\begin{Shaded}
\begin{Highlighting}[]
\CommentTok{\# Sintaxis básica}
\ExtensionTok{adb}\NormalTok{ push archivo\_local ruta\_dispositivo}

\CommentTok{\# Ejemplos:}
\ExtensionTok{adb}\NormalTok{ push documento.pdf /sdcard/Download/}
\ExtensionTok{adb}\NormalTok{ push foto.jpg /sdcard/DCIM/}
\ExtensionTok{adb}\NormalTok{ push video.mp4 /sdcard/Movies/}
\ExtensionTok{adb}\NormalTok{ push musica.mp3 /sdcard/Music/}

\CommentTok{\# Copiar directorio completo}
\ExtensionTok{adb}\NormalTok{ push carpeta/ /sdcard/carpeta/}

\CommentTok{\# Con dispositivo específico}
\ExtensionTok{adb} \AttributeTok{{-}s}\NormalTok{ ABC123 push archivo.txt /sdcard/}
\end{Highlighting}
\end{Shaded}

\subsection{Copiar Archivos del Dispositivo
(Pull)}\label{copiar-archivos-del-dispositivo-pull}

\begin{Shaded}
\begin{Highlighting}[]
\CommentTok{\# Sintaxis básica}
\ExtensionTok{adb}\NormalTok{ pull ruta\_dispositivo archivo\_local}

\CommentTok{\# Ejemplos:}
\ExtensionTok{adb}\NormalTok{ pull /sdcard/Download/documento.pdf ./}
\ExtensionTok{adb}\NormalTok{ pull /sdcard/DCIM/Camera/foto.jpg \textasciitilde{}/Imágenes/}
\ExtensionTok{adb}\NormalTok{ pull /sdcard/WhatsApp/Media/ ./whatsapp\_backup/}

\CommentTok{\# Directorio completo}
\ExtensionTok{adb}\NormalTok{ pull /sdcard/Download/ ./descargas/}

\CommentTok{\# Con dispositivo específico}
\ExtensionTok{adb} \AttributeTok{{-}s}\NormalTok{ ABC123 pull /sdcard/foto.jpg ./}
\end{Highlighting}
\end{Shaded}

\subsection{Explorar Sistema de
Archivos}\label{explorar-sistema-de-archivos}

\begin{Shaded}
\begin{Highlighting}[]
\CommentTok{\# Listar contenido de directorio}
\ExtensionTok{adb}\NormalTok{ shell ls /sdcard/}

\CommentTok{\# Listar con detalles}
\ExtensionTok{adb}\NormalTok{ shell ls }\AttributeTok{{-}la}\NormalTok{ /sdcard/}

\CommentTok{\# Árbol de directorios (si está disponible)}
\ExtensionTok{adb}\NormalTok{ shell find /sdcard/ }\AttributeTok{{-}type}\NormalTok{ d}

\CommentTok{\# Buscar archivos}
\ExtensionTok{adb}\NormalTok{ shell find /sdcard/ }\AttributeTok{{-}name} \StringTok{"*.pdf"}

\CommentTok{\# Ver contenido de archivo}
\ExtensionTok{adb}\NormalTok{ shell cat /sdcard/archivo.txt}

\CommentTok{\# Espacio disponible}
\ExtensionTok{adb}\NormalTok{ shell df }\AttributeTok{{-}h}
\end{Highlighting}
\end{Shaded}

\subsection{Crear/Eliminar Archivos y
Directorios}\label{creareliminar-archivos-y-directorios}

\begin{Shaded}
\begin{Highlighting}[]
\CommentTok{\# Crear directorio}
\ExtensionTok{adb}\NormalTok{ shell mkdir /sdcard/MiCarpeta}

\CommentTok{\# Crear archivo}
\ExtensionTok{adb}\NormalTok{ shell touch /sdcard/archivo.txt}

\CommentTok{\# Escribir en archivo}
\ExtensionTok{adb}\NormalTok{ shell }\StringTok{"echo \textquotesingle{}Hola\textquotesingle{} \textgreater{} /sdcard/archivo.txt"}

\CommentTok{\# Eliminar archivo}
\ExtensionTok{adb}\NormalTok{ shell rm /sdcard/archivo.txt}

\CommentTok{\# Eliminar directorio vacío}
\ExtensionTok{adb}\NormalTok{ shell rmdir /sdcard/MiCarpeta}

\CommentTok{\# Eliminar directorio con contenido}
\ExtensionTok{adb}\NormalTok{ shell rm }\AttributeTok{{-}rf}\NormalTok{ /sdcard/MiCarpeta}
\end{Highlighting}
\end{Shaded}

\begin{center}\rule{0.5\linewidth}{0.5pt}\end{center}

\section{Shell de Android}\label{shell-de-android}

\subsection{Acceder al Shell}\label{acceder-al-shell}

\begin{Shaded}
\begin{Highlighting}[]
\CommentTok{\# Shell interactivo}
\ExtensionTok{adb}\NormalTok{ shell}

\CommentTok{\# Comando único}
\ExtensionTok{adb}\NormalTok{ shell ls /system/bin}

\CommentTok{\# Múltiples comandos}
\ExtensionTok{adb}\NormalTok{ shell }\StringTok{"cd /sdcard \&\& ls {-}l"}

\CommentTok{\# Como root (si disponible)}
\ExtensionTok{adb}\NormalTok{ root}
\ExtensionTok{adb}\NormalTok{ shell}
\CommentTok{\# Ahora estarás como root (\#)}
\end{Highlighting}
\end{Shaded}

\subsection{Comandos Unix Disponibles}\label{comandos-unix-disponibles}

\begin{Shaded}
\begin{Highlighting}[]
\CommentTok{\# Listar herramientas disponibles}
\ExtensionTok{adb}\NormalTok{ shell ls /system/bin}

\CommentTok{\# Ayuda de herramientas}
\ExtensionTok{adb}\NormalTok{ shell toybox }\AttributeTok{{-}{-}help}
\ExtensionTok{adb}\NormalTok{ shell ls }\AttributeTok{{-}{-}help}
\end{Highlighting}
\end{Shaded}

\textbf{Comandos comunes:}

\begin{itemize}
\tightlist
\item
  \texttt{ls} - Listar archivos
\item
  \texttt{cd} - Cambiar directorio
\item
  \texttt{pwd} - Mostrar directorio actual
\item
  \texttt{cat} - Ver contenido de archivo
\item
  \texttt{grep} - Buscar texto
\item
  \texttt{find} - Buscar archivos
\item
  \texttt{ps} - Listar procesos
\item
  \texttt{top} - Monitor de procesos
\item
  \texttt{kill} - Terminar proceso
\item
  \texttt{df} - Espacio en disco
\item
  \texttt{du} - Uso de disco
\item
  \texttt{cp} - Copiar
\item
  \texttt{mv} - Mover
\item
  \texttt{rm} - Eliminar
\item
  \texttt{chmod} - Cambiar permisos
\item
  \texttt{chown} - Cambiar propietario
\end{itemize}

\subsection{Propiedades del Sistema}\label{propiedades-del-sistema}

\begin{Shaded}
\begin{Highlighting}[]
\CommentTok{\# Listar todas las propiedades}
\ExtensionTok{adb}\NormalTok{ shell getprop}

\CommentTok{\# Ver propiedad específica}
\ExtensionTok{adb}\NormalTok{ shell getprop ro.build.version.sdk}

\CommentTok{\# Establecer propiedad (requiere root)}
\ExtensionTok{adb}\NormalTok{ shell setprop debug.myapp.level 5}

\CommentTok{\# Propiedades útiles:}
\ExtensionTok{adb}\NormalTok{ shell getprop ro.product.model           }\CommentTok{\# Modelo}
\ExtensionTok{adb}\NormalTok{ shell getprop ro.build.version.release   }\CommentTok{\# Versión Android}
\ExtensionTok{adb}\NormalTok{ shell getprop ro.build.version.sdk       }\CommentTok{\# Nivel API}
\ExtensionTok{adb}\NormalTok{ shell getprop ro.product.manufacturer    }\CommentTok{\# Fabricante}
\ExtensionTok{adb}\NormalTok{ shell getprop ro.serialno                }\CommentTok{\# Serial}
\ExtensionTok{adb}\NormalTok{ shell getprop net.hostname               }\CommentTok{\# Hostname}
\end{Highlighting}
\end{Shaded}

\subsection{Configuraciones del
Sistema}\label{configuraciones-del-sistema}

\begin{Shaded}
\begin{Highlighting}[]
\CommentTok{\# Ver configuración}
\ExtensionTok{adb}\NormalTok{ shell settings get global wifi\_on}
\ExtensionTok{adb}\NormalTok{ shell settings get system screen\_brightness}

\CommentTok{\# Establecer configuración}
\ExtensionTok{adb}\NormalTok{ shell settings put global airplane\_mode\_on 0}
\ExtensionTok{adb}\NormalTok{ shell settings put system screen\_brightness 100}

\CommentTok{\# Listar todas}
\ExtensionTok{adb}\NormalTok{ shell settings list global}
\ExtensionTok{adb}\NormalTok{ shell settings list system}
\ExtensionTok{adb}\NormalTok{ shell settings list secure}
\end{Highlighting}
\end{Shaded}

\section{Administrador de
Actividades}\label{administrador-de-actividades}

\textbf{El Administrador de Actividades (am) controla actividades,
servicios e intents.}

\subsection{Sintaxis Básica}\label{sintaxis-buxe1sica}

\begin{Shaded}
\begin{Highlighting}[]
\CommentTok{\# Desde shell}
\ExtensionTok{adb}\NormalTok{ shell}
\ExtensionTok{am}\NormalTok{ comando}

\CommentTok{\# Comando directo}
\ExtensionTok{adb}\NormalTok{ shell am comando}
\end{Highlighting}
\end{Shaded}

\subsection{Iniciar Actividad}\label{iniciar-actividad}

\begin{Shaded}
\begin{Highlighting}[]
\CommentTok{\# Iniciar actividad por acción}
\ExtensionTok{adb}\NormalTok{ shell am start }\AttributeTok{{-}a}\NormalTok{ android.intent.action.VIEW}

\CommentTok{\# Iniciar actividad por componente}
\ExtensionTok{adb}\NormalTok{ shell am start }\AttributeTok{{-}n}\NormalTok{ com.android.chrome/com.google.android.apps.chrome.Main}

\CommentTok{\# Abrir URL}
\ExtensionTok{adb}\NormalTok{ shell am start }\AttributeTok{{-}a}\NormalTok{ android.intent.action.VIEW }\AttributeTok{{-}d}\NormalTok{ https://google.com}

\CommentTok{\# Abrir archivo}
\ExtensionTok{adb}\NormalTok{ shell am start }\AttributeTok{{-}a}\NormalTok{ android.intent.action.VIEW }\AttributeTok{{-}d}\NormalTok{ file:///sdcard/documento.pdf}

\CommentTok{\# Con espera hasta completar}
\ExtensionTok{adb}\NormalTok{ shell am start }\AttributeTok{{-}W} \AttributeTok{{-}n}\NormalTok{ com.android.chrome/com.google.android.apps.chrome.Main}

\CommentTok{\# Habilitar debugging}
\ExtensionTok{adb}\NormalTok{ shell am start }\AttributeTok{{-}D} \AttributeTok{{-}n}\NormalTok{ com.example.app/.MainActivity}
\end{Highlighting}
\end{Shaded}

\subsection{Iniciar Servicio}\label{iniciar-servicio}

\begin{Shaded}
\begin{Highlighting}[]
\ExtensionTok{adb}\NormalTok{ shell am startservice }\AttributeTok{{-}n}\NormalTok{ com.example.app/.MyService}
\end{Highlighting}
\end{Shaded}

\subsection{Enviar Broadcast}\label{enviar-broadcast}

\begin{Shaded}
\begin{Highlighting}[]
\CommentTok{\# Broadcast simple}
\ExtensionTok{adb}\NormalTok{ shell am broadcast }\AttributeTok{{-}a}\NormalTok{ android.intent.action.BOOT\_COMPLETED}

\CommentTok{\# Con datos extra}
\ExtensionTok{adb}\NormalTok{ shell am broadcast }\AttributeTok{{-}a}\NormalTok{ com.example.ACTION }\AttributeTok{{-}{-}es}\NormalTok{ key value}
\end{Highlighting}
\end{Shaded}

\subsection{Forzar Detención}\label{forzar-detenciuxf3n}

\begin{Shaded}
\begin{Highlighting}[]
\ExtensionTok{adb}\NormalTok{ shell am force{-}stop com.android.chrome}
\end{Highlighting}
\end{Shaded}

\subsection{Matar Proceso}\label{matar-proceso}

\begin{Shaded}
\begin{Highlighting}[]
\CommentTok{\# Matar app}
\ExtensionTok{adb}\NormalTok{ shell am kill com.android.chrome}

\CommentTok{\# Matar todos los procesos en segundo plano}
\ExtensionTok{adb}\NormalTok{ shell am kill{-}all}
\end{Highlighting}
\end{Shaded}

\subsection{Modificar Pantalla}\label{modificar-pantalla}

\begin{Shaded}
\begin{Highlighting}[]
\CommentTok{\# Cambiar tamaño de pantalla}
\ExtensionTok{adb}\NormalTok{ shell am display{-}size 1920x1080}

\CommentTok{\# Restablecer}
\ExtensionTok{adb}\NormalTok{ shell am display{-}size reset}

\CommentTok{\# Cambiar densidad (DPI)}
\ExtensionTok{adb}\NormalTok{ shell am display{-}density 480}

\CommentTok{\# Restablecer}
\ExtensionTok{adb}\NormalTok{ shell am display{-}density reset}
\end{Highlighting}
\end{Shaded}

\subsection{Ejemplos de Intents}\label{ejemplos-de-intents}

\begin{Shaded}
\begin{Highlighting}[]
\CommentTok{\# Abrir configuración WiFi}
\ExtensionTok{adb}\NormalTok{ shell am start }\AttributeTok{{-}a}\NormalTok{ android.settings.WIFI\_SETTINGS}

\CommentTok{\# Abrir ajustes del sistema}
\ExtensionTok{adb}\NormalTok{ shell am start }\AttributeTok{{-}a}\NormalTok{ android.settings.SETTINGS}

\CommentTok{\# Abrir configuración de aplicación específica}
\ExtensionTok{adb}\NormalTok{ shell am start }\AttributeTok{{-}a}\NormalTok{ android.settings.APPLICATION\_DETAILS\_SETTINGS }\AttributeTok{{-}d}\NormalTok{ package:com.example.app}

\CommentTok{\# Hacer llamada}
\ExtensionTok{adb}\NormalTok{ shell am start }\AttributeTok{{-}a}\NormalTok{ android.intent.action.CALL }\AttributeTok{{-}d}\NormalTok{ tel:123456789}

\CommentTok{\# Enviar SMS}
\ExtensionTok{adb}\NormalTok{ shell am start }\AttributeTok{{-}a}\NormalTok{ android.intent.action.SENDTO }\AttributeTok{{-}d}\NormalTok{ sms:123456789 }\AttributeTok{{-}{-}es}\NormalTok{ sms\_body }\StringTok{"Hola"}

\CommentTok{\# Abrir cámara}
\ExtensionTok{adb}\NormalTok{ shell am start }\AttributeTok{{-}a}\NormalTok{ android.media.action.IMAGE\_CAPTURE}

\CommentTok{\# Compartir texto}
\ExtensionTok{adb}\NormalTok{ shell am start }\AttributeTok{{-}a}\NormalTok{ android.intent.action.SEND }\AttributeTok{{-}{-}es}\NormalTok{ android.intent.extra.TEXT }\StringTok{"Texto"} \AttributeTok{{-}{-}et}\NormalTok{ android.intent.extra.SUBJECT }\StringTok{"Asunto"}
\end{Highlighting}
\end{Shaded}

\section{Administrador de Paquetes}\label{administrador-de-paquetes}

\textbf{El Administrador de Paquetes (pm) gestiona aplicaciones
instaladas.}

\subsection{Listar Paquetes}\label{listar-paquetes}

\begin{Shaded}
\begin{Highlighting}[]
\CommentTok{\# Todos los paquetes}
\ExtensionTok{adb}\NormalTok{ shell pm list packages}

\CommentTok{\# Solo terceros}
\ExtensionTok{adb}\NormalTok{ shell pm list packages }\AttributeTok{{-}3}

\CommentTok{\# Solo sistema}
\ExtensionTok{adb}\NormalTok{ shell pm list packages }\AttributeTok{{-}s}

\CommentTok{\# Habilitados}
\ExtensionTok{adb}\NormalTok{ shell pm list packages }\AttributeTok{{-}e}

\CommentTok{\# Deshabilitados}
\ExtensionTok{adb}\NormalTok{ shell pm list packages }\AttributeTok{{-}d}

\CommentTok{\# Con archivo APK}
\ExtensionTok{adb}\NormalTok{ shell pm list packages }\AttributeTok{{-}f}

\CommentTok{\# Buscar específico}
\ExtensionTok{adb}\NormalTok{ shell pm list packages }\KeywordTok{|} \FunctionTok{grep}\NormalTok{ chrome}
\end{Highlighting}
\end{Shaded}

\subsection{Información del Paquete}\label{informaciuxf3n-del-paquete}

\begin{Shaded}
\begin{Highlighting}[]
\CommentTok{\# Ruta del APK}
\ExtensionTok{adb}\NormalTok{ shell pm path com.android.chrome}

\CommentTok{\# Información completa}
\ExtensionTok{adb}\NormalTok{ shell dumpsys package com.android.chrome}

\CommentTok{\# Permisos}
\ExtensionTok{adb}\NormalTok{ shell pm list permissions}

\CommentTok{\# Permisos de paquete}
\ExtensionTok{adb}\NormalTok{ shell dumpsys package com.android.chrome }\KeywordTok{|} \FunctionTok{grep}\NormalTok{ permission}

\CommentTok{\# Grupos de permisos}
\ExtensionTok{adb}\NormalTok{ shell pm list permission{-}groups}
\end{Highlighting}
\end{Shaded}

\subsection{Otorgar/Revocar Permisos}\label{otorgarrevocar-permisos}

\begin{Shaded}
\begin{Highlighting}[]
\CommentTok{\# Otorgar permiso}
\ExtensionTok{adb}\NormalTok{ shell pm grant com.example.app android.permission.CAMERA}

\CommentTok{\# Revocar permiso}
\ExtensionTok{adb}\NormalTok{ shell pm revoke com.example.app android.permission.CAMERA}

\CommentTok{\# Permisos comunes:}
\CommentTok{\# android.permission.CAMERA}
\CommentTok{\# android.permission.READ\_EXTERNAL\_STORAGE}
\CommentTok{\# android.permission.WRITE\_EXTERNAL\_STORAGE}
\CommentTok{\# android.permission.ACCESS\_FINE\_LOCATION}
\CommentTok{\# android.permission.RECORD\_AUDIO}
\end{Highlighting}
\end{Shaded}

\subsection{Habilitar/Deshabilitar
Paquete}\label{habilitardeshabilitar-paquete}

\begin{Shaded}
\begin{Highlighting}[]
\CommentTok{\# Deshabilitar}
\ExtensionTok{adb}\NormalTok{ shell pm disable{-}user com.example.app}

\CommentTok{\# Habilitar}
\ExtensionTok{adb}\NormalTok{ shell pm enable com.example.app}
\end{Highlighting}
\end{Shaded}

\subsection{Limpiar Datos}\label{limpiar-datos}

\begin{Shaded}
\begin{Highlighting}[]
\CommentTok{\# Borrar datos de aplicación}
\ExtensionTok{adb}\NormalTok{ shell pm clear com.android.chrome}
\end{Highlighting}
\end{Shaded}

\subsection{Usuarios}\label{usuarios}

\begin{Shaded}
\begin{Highlighting}[]
\CommentTok{\# Listar usuarios}
\ExtensionTok{adb}\NormalTok{ shell pm list users}

\CommentTok{\# Crear usuario}
\ExtensionTok{adb}\NormalTok{ shell pm create{-}user }\StringTok{"Nombre Usuario"}

\CommentTok{\# Eliminar usuario}
\ExtensionTok{adb}\NormalTok{ shell pm remove{-}user USER\_ID}

\CommentTok{\# Máximo de usuarios}
\ExtensionTok{adb}\NormalTok{ shell pm get{-}max{-}users}
\end{Highlighting}
\end{Shaded}

\section{Captura de Pantalla y
Grabación}\label{captura-de-pantalla-y-grabaciuxf3n}

\subsection{Captura de Pantalla}\label{captura-de-pantalla}

\begin{Shaded}
\begin{Highlighting}[]
\CommentTok{\# Capturar y guardar en dispositivo}
\ExtensionTok{adb}\NormalTok{ shell screencap /sdcard/screenshot.png}

\CommentTok{\# Capturar y descargar}
\ExtensionTok{adb}\NormalTok{ shell screencap /sdcard/screenshot.png}
\ExtensionTok{adb}\NormalTok{ pull /sdcard/screenshot.png ./}
\ExtensionTok{adb}\NormalTok{ shell rm /sdcard/screenshot.png}

\CommentTok{\# Captura directa (un solo comando)}
\ExtensionTok{adb}\NormalTok{ exec{-}out screencap }\AttributeTok{{-}p} \OperatorTok{\textgreater{}}\NormalTok{ screenshot.png}

\CommentTok{\# Con timestamp}
\ExtensionTok{adb}\NormalTok{ exec{-}out screencap }\AttributeTok{{-}p} \OperatorTok{\textgreater{}}\NormalTok{ screenshot\_}\VariableTok{$(}\FunctionTok{date}\NormalTok{ +\%Y\%m\%d\_\%H\%M\%S}\VariableTok{)}\NormalTok{.png}
\end{Highlighting}
\end{Shaded}

\subsection{Grabación de Video}\label{grabaciuxf3n-de-video}

\begin{Shaded}
\begin{Highlighting}[]
\CommentTok{\# Grabar video (hasta 3 minutos por defecto)}
\ExtensionTok{adb}\NormalTok{ shell screenrecord /sdcard/video.mp4}

\CommentTok{\# Detener con Ctrl+C}

\CommentTok{\# Descargar}
\ExtensionTok{adb}\NormalTok{ pull /sdcard/video.mp4 ./}

\CommentTok{\# Con límite de tiempo (segundos)}
\ExtensionTok{adb}\NormalTok{ shell screenrecord }\AttributeTok{{-}{-}time{-}limit}\NormalTok{ 30 /sdcard/video.mp4}

\CommentTok{\# Especificar tamaño}
\ExtensionTok{adb}\NormalTok{ shell screenrecord }\AttributeTok{{-}{-}size}\NormalTok{ 1280x720 /sdcard/video.mp4}

\CommentTok{\# Cambiar bitrate (por defecto 20 Mbps)}
\ExtensionTok{adb}\NormalTok{ shell screenrecord }\AttributeTok{{-}{-}bit{-}rate}\NormalTok{ 8000000 /sdcard/video.mp4}

\CommentTok{\# Verbose (mostrar info)}
\ExtensionTok{adb}\NormalTok{ shell screenrecord }\AttributeTok{{-}{-}verbose}\NormalTok{ /sdcard/video.mp4}

\CommentTok{\# Rotar video 90 grados}
\ExtensionTok{adb}\NormalTok{ shell screenrecord }\AttributeTok{{-}{-}rotate}\NormalTok{ /sdcard/video.mp4}
\end{Highlighting}
\end{Shaded}

\textbf{Limitaciones de screenrecord:}

\begin{itemize}
\tightlist
\item
  No graba audio
\item
  Máximo 3 minutos por video
\item
  No funciona en Wear OS
\item
  No soporta rotación durante grabación
\end{itemize}

\section{Redirección de Puertos}\label{redirecciuxf3n-de-puertos}

\subsection{Forward (Host →
Dispositivo)}\label{forward-host-dispositivo}

\begin{Shaded}
\begin{Highlighting}[]
\CommentTok{\# Redirigir puerto host 6100 a dispositivo 7100}
\ExtensionTok{adb}\NormalTok{ forward tcp:6100 tcp:7100}

\CommentTok{\# Listar redirecciones}
\ExtensionTok{adb}\NormalTok{ forward }\AttributeTok{{-}{-}list}

\CommentTok{\# Eliminar redirección específica}
\ExtensionTok{adb}\NormalTok{ forward }\AttributeTok{{-}{-}remove}\NormalTok{ tcp:6100}

\CommentTok{\# Eliminar todas}
\ExtensionTok{adb}\NormalTok{ forward }\AttributeTok{{-}{-}remove{-}all}

\CommentTok{\# Ejemplo: acceder a servidor en dispositivo desde host}
\CommentTok{\# Si hay servidor en dispositivo:puerto 8080}
\ExtensionTok{adb}\NormalTok{ forward tcp:8080 tcp:8080}
\CommentTok{\# Ahora accesible en: http://localhost:8080}
\end{Highlighting}
\end{Shaded}

\subsection{Reverse (Dispositivo →
Host)}\label{reverse-dispositivo-host}

\begin{Shaded}
\begin{Highlighting}[]
\CommentTok{\# Redirigir puerto dispositivo 9000 a host 9000}
\ExtensionTok{adb}\NormalTok{ reverse tcp:9000 tcp:9000}

\CommentTok{\# Listar}
\ExtensionTok{adb}\NormalTok{ reverse }\AttributeTok{{-}{-}list}

\CommentTok{\# Eliminar}
\ExtensionTok{adb}\NormalTok{ reverse }\AttributeTok{{-}{-}remove}\NormalTok{ tcp:9000}

\CommentTok{\# Eliminar todas}
\ExtensionTok{adb}\NormalTok{ reverse }\AttributeTok{{-}{-}remove{-}all}

\CommentTok{\# Ejemplo: dispositivo accede a servidor en host}
\CommentTok{\# Si hay servidor en host:puerto 3000}
\ExtensionTok{adb}\NormalTok{ reverse tcp:3000 tcp:3000}
\CommentTok{\# Dispositivo puede acceder a http://localhost:3000}
\end{Highlighting}
\end{Shaded}

\section{Comandos ADB para Conectar y Instalar en
Windows}\label{comandos-adb-para-conectar-y-instalar-en-windows}

\subsection{Versión de Microsoft Windows:
10.0.22000.282}\label{versiuxf3n-de-microsoft-windows-10.0.22000.282}

\begin{enumerate}
\def\labelenumi{(\alph{enumi})}
\setcounter{enumi}{2}
\tightlist
\item
  Microsoft Corporation. Todos los derechos reservados.
\end{enumerate}

\subsection{\texorpdfstring{Cambio de directorio a la ubicación de
\texttt{ADB}}{Cambio de directorio a la ubicación de ADB}}\label{cambio-de-directorio-a-la-ubicaciuxf3n-de-adb}

\begin{Shaded}
\begin{Highlighting}[]

\ExtensionTok{C:\textbackslash{}Users\textbackslash{}achalmaedison}\OperatorTok{\textgreater{}}\NormalTok{cd }\StringTok{"C:\textbackslash{}Users\textbackslash{}achalmaedison\textbackslash{}Downloads\textbackslash{}ADB"}
\end{Highlighting}
\end{Shaded}

\subsection{Conexión con ADB}\label{conexiuxf3n-con-adb}

\begin{Shaded}
\begin{Highlighting}[]

\ExtensionTok{C:\textbackslash{}Users\textbackslash{}achalmaedison\textbackslash{}Downloads\textbackslash{}ADB}\OperatorTok{\textgreater{}}\NormalTok{adb connect 127.0.0.1:58526}
\end{Highlighting}
\end{Shaded}

\textbf{Salida esperada:}

\begin{Shaded}
\begin{Highlighting}[]

\ExtensionTok{Ya}\NormalTok{ conectado a 127.0.0.1:58526}
\end{Highlighting}
\end{Shaded}

\subsection{Instalación de Instagram
APK}\label{instalaciuxf3n-de-instagram-apk}

\begin{Shaded}
\begin{Highlighting}[]

\ExtensionTok{C:\textbackslash{}Users\textbackslash{}achalmaedison\textbackslash{}Downloads\textbackslash{}ADB}\OperatorTok{\textgreater{}}\NormalTok{adb install Instagram.apk}
\end{Highlighting}
\end{Shaded}

\textbf{Salida esperada:}

\begin{Shaded}
\begin{Highlighting}[]

\ExtensionTok{Realizando}\NormalTok{ Instalación en Streaming}

\ExtensionTok{Instalación}\NormalTok{ Exitosa}
\end{Highlighting}
\end{Shaded}

Este es un ejemplo de cómo utilizar los comandos ADB en un entorno
Windows.

\subsection{Descarga de ADB}\label{descarga-de-adb}

Puedes descargar ADB desde el siguiente enlace:

\href{https://developer.android.com/studio/releases/platform-tools}{Plataforma
de herramientas de Android}

\subsection{Configuración}\label{configuraciuxf3n}

\begin{enumerate}
\def\labelenumi{\arabic{enumi}.}
\item
  Configurar el subsistema de Android.
\item
  Activar el modo desarrollador.
\item
  Copiar la dirección IP del dispositivo.
\item
  Abrir \textbf{CMD} como administrador.
\item
  Navegar a la carpeta \texttt{platform-tools} ubicada en
  \texttt{C:\textbackslash{}platform-tools}:
\end{enumerate}

\begin{Shaded}
\begin{Highlighting}[]

\BuiltInTok{cd}\NormalTok{ C:}\DataTypeTok{\textbackslash{}p}\NormalTok{latform{-}tools}
\end{Highlighting}
\end{Shaded}

\subsection{Comandos para conexión e instalación de
aplicaciones}\label{comandos-para-conexiuxf3n-e-instalaciuxf3n-de-aplicaciones}

Conectar el dispositivo:

\begin{Shaded}
\begin{Highlighting}[]

\ExtensionTok{adb}\NormalTok{ connect 127.0.0.1:58526}
\end{Highlighting}
\end{Shaded}

Instalar aplicaciones:

\begin{Shaded}
\begin{Highlighting}[]

\ExtensionTok{adb}\NormalTok{ install gmail.apk}
\end{Highlighting}
\end{Shaded}

\begin{Shaded}
\begin{Highlighting}[]

\ExtensionTok{adb}\NormalTok{ install facebook{-}354{-}0{-}0{-}22{-}110.apk}
\end{Highlighting}
\end{Shaded}

\textbf{Nota:} Si los comandos no funcionan, intenta reiniciar el
sistema.

\subsection{Ejemplo de salida
esperada}\label{ejemplo-de-salida-esperada}

\begin{Shaded}
\begin{Highlighting}[]

\ExtensionTok{Microsoft}\NormalTok{ Windows [Version 10.0.22000.527]}

\KeywordTok{(}\ExtensionTok{c}\KeywordTok{)} \ExtensionTok{Microsoft}\NormalTok{ Corporation. All rights reserved.}
\end{Highlighting}
\end{Shaded}

\begin{Shaded}
\begin{Highlighting}[]

\ExtensionTok{C:\textbackslash{}platform{-}tools}\OperatorTok{\textgreater{}}\NormalTok{adb connect 172.28.22.78}

\ExtensionTok{*}\NormalTok{ daemon not running}\KeywordTok{;} \ExtensionTok{starting}\NormalTok{ now at tcp:5037}

\ExtensionTok{*}\NormalTok{ daemon started successfully}

\ExtensionTok{connected}\NormalTok{ to 172.28.22.78:5555}
\end{Highlighting}
\end{Shaded}

\begin{Shaded}
\begin{Highlighting}[]

\ExtensionTok{C:\textbackslash{}platform{-}tools}\OperatorTok{\textgreater{}}\NormalTok{adb install termux{-}app\_v0.118.0+github{-}debug\_arm64{-}v8a.apk}

\ExtensionTok{Performing}\NormalTok{ Streamed Install}

\ExtensionTok{Success}
\end{Highlighting}
\end{Shaded}

\begin{Shaded}
\begin{Highlighting}[]

\ExtensionTok{C:\textbackslash{}platform{-}tools}\OperatorTok{\textgreater{}}
\end{Highlighting}
\end{Shaded}

\section{Windows 11: Cómo descargar APK usando el Subsistema de Windows
para Android y
ADB}\label{windows-11-cuxf3mo-descargar-apk-usando-el-subsistema-de-windows-para-android-y-adb}

Aquí te explicamos cómo descargar un archivo APK para instalar una
aplicación de Android en tu PC con Windows 11 usando el Subsistema de
Windows para Android. Puedes instalar el Subsistema de Windows para
Android manualmente en tu PC con su archivo Msixbundle siguiendo nuestra
\href{https://nerdschalk.com/android-apps-on-windows-11-dev-channel-how-to-install-windows-subsystem-for-android-manually-with-msixbundle/}{guía
aquí}.

\subsection{Paso 1: Habilitar el modo de desarrollador en el subsistema
de
Windows}\label{paso-1-habilitar-el-modo-de-desarrollador-en-el-subsistema-de-windows}

\begin{enumerate}
\def\labelenumi{\arabic{enumi}.}
\tightlist
\item
  Instala el
  \href{https://nerdschalk.com/android-apps-on-windows-11-dev-channel-how-to-install-windows-subsystem-for-android-manually-with-msixbundle/}{Subsistema
  de Windows para Android}.
\item
  Abre la aplicación `Subsistema de Windows para Android' en tu PC. Para
  ello, presiona la tecla \textbf{Windows} y busca \textbf{Subsistema de
  Windows para Android}.
\item
  Haz clic en la aplicación para abrirla.
\item
  Dentro de la aplicación, activa el \textbf{Modo Desarrollador}.
\end{enumerate}

\subsection{Paso 2: Instalar las herramientas de la plataforma
SDK}\label{paso-2-instalar-las-herramientas-de-la-plataforma-sdk}

\begin{enumerate}
\def\labelenumi{\arabic{enumi}.}
\tightlist
\item
  Visita la página de herramientas de la plataforma SDK de Google
  \href{https://developer.android.com/studio/releases/platform-tools.html}{aquí}.
\item
  Descarga \textbf{SDK Platform-Tools para Windows}.
\item
  Acepta los términos y condiciones y haz clic en el botón de descarga.
\item
  Se descargará un archivo ZIP llamado
  \textbf{platform-tools\_rXX.X.X-windows.zip} (la versión puede
  variar).
\item
  Crea una carpeta separada en el Explorador de Windows, por ejemplo,
  \texttt{C:\textbackslash{}Plataforma-Tools}.
\item
  Mueve el archivo ZIP descargado a esta carpeta.
\item
  Haz clic derecho en el archivo y selecciona \textbf{Extraer todo},
  luego haz clic en \textbf{Extraer}.
\item
  Abre la carpeta \texttt{platform-tools}, donde encontrarás
  \texttt{adb.exe} y otros archivos.
\end{enumerate}

\subsection{Paso 3: Instalar la aplicación de
Android}\label{paso-3-instalar-la-aplicaciuxf3n-de-android}

\begin{enumerate}
\def\labelenumi{\arabic{enumi}.}
\item
  Abre la carpeta \textbf{platform-tools}.
\item
  Haz clic en la barra de direcciones, escribe \textbf{\texttt{cmd}} y
  presiona \textbf{Enter}.
\item
  Se abrirá una ventana de comandos en la ubicación de la carpeta
  \textbf{platform-tools}.
\item
  Descarga el archivo APK de la aplicación de Android que deseas
  instalar.

  \begin{itemize}
  \tightlist
  \item
    Por ejemplo, para instalar Snapchat, busca \textbf{Snapchat APK} en
    Google y descarga el archivo de una fuente confiable.
  \item
    Renombra el archivo a algo simple, como \texttt{snapchat.apk}, y
    muévelo a la carpeta \textbf{platform-tools}.
  \end{itemize}
\item
  Abre el \textbf{Subsistema de Windows para Android} y copia la
  dirección \textbf{IP} en la opción de \textbf{Modo Desarrollador}.
\item
  En la ventana de comandos, ejecuta el siguiente comando:

\begin{Shaded}
\begin{Highlighting}[]
\ExtensionTok{adb.exe}\NormalTok{ connect }\PreprocessorTok{[}\SpecialStringTok{DIRECCIÓN\_IP}\PreprocessorTok{]}
\end{Highlighting}
\end{Shaded}

  \textbf{Ejemplo:}

\begin{Shaded}
\begin{Highlighting}[]
\ExtensionTok{adb.exe}\NormalTok{ connect 127.0.0.1:12345}
\end{Highlighting}
\end{Shaded}
\item
  Luego, instala la aplicación ejecutando:

\begin{Shaded}
\begin{Highlighting}[]
\ExtensionTok{adb.exe}\NormalTok{ install }\PreprocessorTok{[}\SpecialStringTok{NOMBRE\_DEL\_APK}\PreprocessorTok{]}
\end{Highlighting}
\end{Shaded}

  \textbf{Ejemplo:}

\begin{Shaded}
\begin{Highlighting}[]
\ExtensionTok{adb.exe}\NormalTok{ install snapchat.apk}
\end{Highlighting}
\end{Shaded}
\item
  Cuando la instalación finalice, verás el mensaje \textbf{Success}.
\item
  Cierra la ventana de comandos.
\item
  Abre la aplicación en tu PC escribiendo su nombre en el menú Inicio
  (por ejemplo, \textbf{Snapchat}).
\end{enumerate}

\section{Casos de Uso Profesionales}\label{casos-de-uso-profesionales}

\subsection{Caso 1: Instalación Masiva de
APKs}\label{caso-1-instalaciuxf3n-masiva-de-apks}

\textbf{Escenario:} Instalar múltiples APKs en varios dispositivos.

\begin{Shaded}
\begin{Highlighting}[]
\CommentTok{\#!/bin/bash}
\CommentTok{\# instalar\_apks.sh}

\CommentTok{\# Directorio con APKs}
\VariableTok{APK\_DIR}\OperatorTok{=}\StringTok{"./apks"}

\CommentTok{\# Obtener lista de dispositivos}
\VariableTok{devices}\OperatorTok{=}\VariableTok{$(}\ExtensionTok{adb}\NormalTok{ devices }\KeywordTok{|} \FunctionTok{grep} \AttributeTok{{-}v} \StringTok{"List"} \KeywordTok{|} \FunctionTok{grep} \StringTok{"device$"} \KeywordTok{|} \FunctionTok{awk} \StringTok{\textquotesingle{}\{print $1\}\textquotesingle{}}\VariableTok{)}

\CommentTok{\# Para cada dispositivo}
\ControlFlowTok{for}\NormalTok{ device }\KeywordTok{in} \VariableTok{$devices}\KeywordTok{;} \ControlFlowTok{do}
    \BuiltInTok{echo} \StringTok{"Instalando en dispositivo: }\VariableTok{$device}\StringTok{"}
    
    \CommentTok{\# Para cada APK}
    \ControlFlowTok{for}\NormalTok{ apk }\KeywordTok{in} \StringTok{"}\VariableTok{$APK\_DIR}\StringTok{"}\NormalTok{/}\PreprocessorTok{*}\NormalTok{.apk}\KeywordTok{;} \ControlFlowTok{do}
        \BuiltInTok{echo} \StringTok{"  Instalando: }\VariableTok{$(}\FunctionTok{basename} \VariableTok{$apk)}\StringTok{"}
        \ExtensionTok{adb} \AttributeTok{{-}s} \StringTok{"}\VariableTok{$device}\StringTok{"}\NormalTok{ install }\AttributeTok{{-}r} \StringTok{"}\VariableTok{$apk}\StringTok{"}
    \ControlFlowTok{done}
\ControlFlowTok{done}

\BuiltInTok{echo} \StringTok{"Instalación completada"}
\end{Highlighting}
\end{Shaded}

\textbf{Usar:}

\begin{Shaded}
\begin{Highlighting}[]
\FunctionTok{chmod}\NormalTok{ +x instalar\_apks.sh}
\ExtensionTok{./instalar\_apks.sh}
\end{Highlighting}
\end{Shaded}

\subsection{Caso 2: Backup Automático de
Aplicación}\label{caso-2-backup-automuxe1tico-de-aplicaciuxf3n}

\textbf{Escenario:} Respaldar datos de aplicación específica.

\begin{Shaded}
\begin{Highlighting}[]
\CommentTok{\#!/bin/bash}
\CommentTok{\# backup\_app.sh}

\VariableTok{PACKAGE}\OperatorTok{=}\StringTok{"com.example.app"}
\VariableTok{BACKUP\_DIR}\OperatorTok{=}\StringTok{"./backups/}\VariableTok{$(}\FunctionTok{date}\NormalTok{ +\%Y\%m\%d\_\%H\%M\%S}\VariableTok{)}\StringTok{"}

\FunctionTok{mkdir} \AttributeTok{{-}p} \StringTok{"}\VariableTok{$BACKUP\_DIR}\StringTok{"}

\BuiltInTok{echo} \StringTok{"Respaldando }\VariableTok{$PACKAGE}\StringTok{..."}

\CommentTok{\# Backup de APK}
\BuiltInTok{echo} \StringTok{"1. Respaldando APK..."}
\VariableTok{APK\_PATH}\OperatorTok{=}\VariableTok{$(}\ExtensionTok{adb}\NormalTok{ shell pm path }\StringTok{"}\VariableTok{$PACKAGE}\StringTok{"} \KeywordTok{|} \FunctionTok{cut} \AttributeTok{{-}d:} \AttributeTok{{-}f2} \KeywordTok{|} \FunctionTok{tr} \AttributeTok{{-}d} \StringTok{\textquotesingle{}\textbackslash{}r\textquotesingle{}}\VariableTok{)}
\ExtensionTok{adb}\NormalTok{ pull }\StringTok{"}\VariableTok{$APK\_PATH}\StringTok{"} \StringTok{"}\VariableTok{$BACKUP\_DIR}\StringTok{/app.apk"}

\CommentTok{\# Backup de datos (requiere app debuggable o root)}
\BuiltInTok{echo} \StringTok{"2. Respaldando datos..."}
\ExtensionTok{adb}\NormalTok{ backup }\AttributeTok{{-}f} \StringTok{"}\VariableTok{$BACKUP\_DIR}\StringTok{/data.ab"} \AttributeTok{{-}noapk} \StringTok{"}\VariableTok{$PACKAGE}\StringTok{"}

\CommentTok{\# Backup de base de datos (si accesible)}
\BuiltInTok{echo} \StringTok{"3. Respaldando bases de datos..."}
\ExtensionTok{adb}\NormalTok{ pull /data/data/}\StringTok{"}\VariableTok{$PACKAGE}\StringTok{"}\NormalTok{/databases }\StringTok{"}\VariableTok{$BACKUP\_DIR}\StringTok{/databases/"}

\BuiltInTok{echo} \StringTok{"Backup completado en: }\VariableTok{$BACKUP\_DIR}\StringTok{"}
\end{Highlighting}
\end{Shaded}

\subsection{Caso 3: Testing
Automatizado}\label{caso-3-testing-automatizado}

\textbf{Escenario:} Ejecutar suite de tests en dispositivo.

\begin{Shaded}
\begin{Highlighting}[]
\CommentTok{\#!/bin/bash}
\CommentTok{\# run\_tests.sh}

\VariableTok{PACKAGE}\OperatorTok{=}\StringTok{"com.example.app"}
\VariableTok{TEST\_PACKAGE}\OperatorTok{=}\StringTok{"com.example.app.test"}

\BuiltInTok{echo} \StringTok{"Iniciando tests..."}

\CommentTok{\# Instalar APKs}
\ExtensionTok{adb}\NormalTok{ install }\AttributeTok{{-}r}\NormalTok{ app.apk}
\ExtensionTok{adb}\NormalTok{ install }\AttributeTok{{-}r}\NormalTok{ app{-}test.apk}

\CommentTok{\# Ejecutar tests}
\ExtensionTok{adb}\NormalTok{ shell am instrument }\AttributeTok{{-}w} \DataTypeTok{\textbackslash{}}
    \AttributeTok{{-}e}\NormalTok{ package }\StringTok{"}\VariableTok{$PACKAGE}\StringTok{"} \DataTypeTok{\textbackslash{}}
    \StringTok{"}\VariableTok{$TEST\_PACKAGE}\StringTok{/androidx.test.runner.AndroidJUnitRunner"}

\CommentTok{\# Generar reporte}
\ExtensionTok{adb}\NormalTok{ pull /sdcard/test{-}results.xml ./}
\end{Highlighting}
\end{Shaded}

\subsection{Caso 4: Análisis de
Rendimiento}\label{caso-4-anuxe1lisis-de-rendimiento}

\textbf{Escenario:} Capturar métricas de rendimiento de aplicación.

\begin{Shaded}
\begin{Highlighting}[]
\CommentTok{\#!/bin/bash}
\CommentTok{\# performance\_test.sh}

\VariableTok{PACKAGE}\OperatorTok{=}\StringTok{"com.example.app"}
\VariableTok{DURATION}\OperatorTok{=}\NormalTok{60  }\CommentTok{\# segundos}

\BuiltInTok{echo} \StringTok{"Analizando rendimiento de }\VariableTok{$PACKAGE}\StringTok{ por }\VariableTok{$DURATION}\StringTok{ segundos..."}

\CommentTok{\# Limpiar logs}
\ExtensionTok{adb}\NormalTok{ logcat }\AttributeTok{{-}c}

\CommentTok{\# Iniciar app}
\ExtensionTok{adb}\NormalTok{ shell monkey }\AttributeTok{{-}p} \StringTok{"}\VariableTok{$PACKAGE}\StringTok{"} \AttributeTok{{-}c}\NormalTok{ android.intent.category.LAUNCHER 1}

\CommentTok{\# Capturar CPU/Memoria}
\ExtensionTok{adb}\NormalTok{ shell top }\AttributeTok{{-}n}\NormalTok{ 1 }\AttributeTok{{-}d}\NormalTok{ 1 }\AttributeTok{{-}m}\NormalTok{ 10 }\OperatorTok{\textgreater{}}\NormalTok{ cpu\_usage.txt }\KeywordTok{\&}
\VariableTok{TOP\_PID}\OperatorTok{=}\VariableTok{$!}

\CommentTok{\# Esperar duración}
\FunctionTok{sleep} \StringTok{"}\VariableTok{$DURATION}\StringTok{"}

\CommentTok{\# Detener captura}
\BuiltInTok{kill} \VariableTok{$TOP\_PID}

\CommentTok{\# Obtener memoria detallada}
\ExtensionTok{adb}\NormalTok{ shell dumpsys meminfo }\StringTok{"}\VariableTok{$PACKAGE}\StringTok{"} \OperatorTok{\textgreater{}}\NormalTok{ memory\_info.txt}

\CommentTok{\# Obtener logs}
\ExtensionTok{adb}\NormalTok{ logcat }\AttributeTok{{-}d} \OperatorTok{\textgreater{}}\NormalTok{ app\_logs.txt}

\BuiltInTok{echo} \StringTok{"Análisis completado"}
\end{Highlighting}
\end{Shaded}

\subsection{Caso 5: Captura Automatizada de
Screenshots}\label{caso-5-captura-automatizada-de-screenshots}

\textbf{Escenario:} Capturar screenshots para documentación.

\begin{Shaded}
\begin{Highlighting}[]
\CommentTok{\#!/bin/bash}
\CommentTok{\# capture\_screenshots.sh}

\VariableTok{OUTPUT\_DIR}\OperatorTok{=}\StringTok{"./screenshots/}\VariableTok{$(}\FunctionTok{date}\NormalTok{ +\%Y\%m\%d\_\%H\%M\%S}\VariableTok{)}\StringTok{"}
\FunctionTok{mkdir} \AttributeTok{{-}p} \StringTok{"}\VariableTok{$OUTPUT\_DIR}\StringTok{"}

\VariableTok{SCREENS}\OperatorTok{=}\VariableTok{(}
    \StringTok{"com.example.app/.MainActivity:Pantalla Principal"}
    \StringTok{"com.example.app/.SettingsActivity:Configuración"}
    \StringTok{"com.example.app/.ProfileActivity:Perfil"}
\VariableTok{)}

\ControlFlowTok{for}\NormalTok{ screen }\KeywordTok{in} \StringTok{"}\VariableTok{$\{SCREENS}\OperatorTok{[@]}\VariableTok{\}}\StringTok{"}\KeywordTok{;} \ControlFlowTok{do}
    \VariableTok{ACTIVITY}\OperatorTok{=}\VariableTok{$(}\BuiltInTok{echo} \StringTok{"}\VariableTok{$screen}\StringTok{"} \KeywordTok{|} \FunctionTok{cut} \AttributeTok{{-}d:} \AttributeTok{{-}f1}\VariableTok{)}
    \VariableTok{NAME}\OperatorTok{=}\VariableTok{$(}\BuiltInTok{echo} \StringTok{"}\VariableTok{$screen}\StringTok{"} \KeywordTok{|} \FunctionTok{cut} \AttributeTok{{-}d:} \AttributeTok{{-}f2}\VariableTok{)}
    
    \BuiltInTok{echo} \StringTok{"Capturando: }\VariableTok{$NAME}\StringTok{"}
    
    \CommentTok{\# Abrir actividad}
    \ExtensionTok{adb}\NormalTok{ shell am start }\AttributeTok{{-}n} \StringTok{"}\VariableTok{$ACTIVITY}\StringTok{"}
    \FunctionTok{sleep}\NormalTok{ 2}
    
    \CommentTok{\# Capturar}
    \VariableTok{FILENAME}\OperatorTok{=}\VariableTok{$(}\BuiltInTok{echo} \StringTok{"}\VariableTok{$NAME}\StringTok{"} \KeywordTok{|} \FunctionTok{tr} \StringTok{\textquotesingle{} \textquotesingle{}} \StringTok{\textquotesingle{}\_\textquotesingle{}} \KeywordTok{|} \FunctionTok{tr} \StringTok{\textquotesingle{}[:upper:]\textquotesingle{}} \StringTok{\textquotesingle{}[:lower:]\textquotesingle{}}\VariableTok{)}
    \ExtensionTok{adb}\NormalTok{ exec{-}out screencap }\AttributeTok{{-}p} \OperatorTok{\textgreater{}} \StringTok{"}\VariableTok{$OUTPUT\_DIR}\StringTok{/}\VariableTok{$\{FILENAME\}}\StringTok{.png"}
    
    \FunctionTok{sleep}\NormalTok{ 1}
\ControlFlowTok{done}

\BuiltInTok{echo} \StringTok{"Screenshots guardados en: }\VariableTok{$OUTPUT\_DIR}\StringTok{"}
\end{Highlighting}
\end{Shaded}

\subsection{Caso 6: Monitoreo de Logs en Tiempo
Real}\label{caso-6-monitoreo-de-logs-en-tiempo-real}

\textbf{Escenario:} Filtrar logs de aplicación específica.

\begin{Shaded}
\begin{Highlighting}[]
\CommentTok{\#!/bin/bash}
\CommentTok{\# monitor\_logs.sh}

\VariableTok{PACKAGE}\OperatorTok{=}\StringTok{"com.example.app"}

\CommentTok{\# Obtener PID}
\VariableTok{PID}\OperatorTok{=}\VariableTok{$(}\ExtensionTok{adb}\NormalTok{ shell ps }\KeywordTok{|} \FunctionTok{grep} \StringTok{"}\VariableTok{$PACKAGE}\StringTok{"} \KeywordTok{|} \FunctionTok{awk} \StringTok{\textquotesingle{}\{print $2\}\textquotesingle{}}\VariableTok{)}

\BuiltInTok{echo} \StringTok{"Monitoreando logs de }\VariableTok{$PACKAGE}\StringTok{ (PID: }\VariableTok{$PID}\StringTok{)"}
\BuiltInTok{echo} \StringTok{"Presiona Ctrl+C para detener"}

\CommentTok{\# Filtrar logs por PID}
\ExtensionTok{adb}\NormalTok{ logcat }\AttributeTok{{-}{-}pid}\OperatorTok{=}\StringTok{"}\VariableTok{$PID}\StringTok{"} \AttributeTok{{-}v}\NormalTok{ time}
\end{Highlighting}
\end{Shaded}

\subsection{Caso 7: Instalación y Testing
Rápido}\label{caso-7-instalaciuxf3n-y-testing-ruxe1pido}

\textbf{Escenario:} Ciclo desarrollo-testing rápido.

\begin{Shaded}
\begin{Highlighting}[]
\CommentTok{\#!/bin/bash}
\CommentTok{\# quick\_test.sh}

\VariableTok{APK}\OperatorTok{=}\StringTok{"app{-}debug.apk"}
\VariableTok{PACKAGE}\OperatorTok{=}\StringTok{"com.example.app"}

\BuiltInTok{echo} \StringTok{"1. Desinstalando versión anterior..."}
\ExtensionTok{adb}\NormalTok{ uninstall }\StringTok{"}\VariableTok{$PACKAGE}\StringTok{"} \DecValTok{2}\OperatorTok{\textgreater{}}\NormalTok{/dev/null}

\BuiltInTok{echo} \StringTok{"2. Instalando nueva versión..."}
\ExtensionTok{adb}\NormalTok{ install }\StringTok{"}\VariableTok{$APK}\StringTok{"}

\BuiltInTok{echo} \StringTok{"3. Otorgando permisos..."}
\ExtensionTok{adb}\NormalTok{ shell pm grant }\StringTok{"}\VariableTok{$PACKAGE}\StringTok{"}\NormalTok{ android.permission.CAMERA}
\ExtensionTok{adb}\NormalTok{ shell pm grant }\StringTok{"}\VariableTok{$PACKAGE}\StringTok{"}\NormalTok{ android.permission.READ\_EXTERNAL\_STORAGE}
\ExtensionTok{adb}\NormalTok{ shell pm grant }\StringTok{"}\VariableTok{$PACKAGE}\StringTok{"}\NormalTok{ android.permission.WRITE\_EXTERNAL\_STORAGE}

\BuiltInTok{echo} \StringTok{"4. Iniciando app..."}
\ExtensionTok{adb}\NormalTok{ shell monkey }\AttributeTok{{-}p} \StringTok{"}\VariableTok{$PACKAGE}\StringTok{"} \AttributeTok{{-}c}\NormalTok{ android.intent.category.LAUNCHER 1}

\BuiltInTok{echo} \StringTok{"5. Mostrando logs..."}
\ExtensionTok{adb}\NormalTok{ logcat }\AttributeTok{{-}c}
\ExtensionTok{adb}\NormalTok{ logcat }\KeywordTok{|} \FunctionTok{grep} \StringTok{"}\VariableTok{$PACKAGE}\StringTok{"}
\end{Highlighting}
\end{Shaded}

\section{Troubleshooting}\label{troubleshooting}

\subsection{Problema 1: ``adb: command not
found''}\label{problema-1-adb-command-not-found}

\textbf{Causa:} ADB no está instalado o no está en PATH.

\textbf{Soluciones:}

\begin{Shaded}
\begin{Highlighting}[]
\CommentTok{\# 1. Verificar instalación}
\FunctionTok{which}\NormalTok{ adb}

\CommentTok{\# 2. Si no está, instalar}
\FunctionTok{sudo}\NormalTok{ pacman }\AttributeTok{{-}S}\NormalTok{ android{-}tools}

\CommentTok{\# 3. Verificar PATH}
\BuiltInTok{echo} \VariableTok{$PATH} \KeywordTok{|} \FunctionTok{grep}\NormalTok{ android}

\CommentTok{\# 4. Si instalaste manualmente, agregar a PATH}
\BuiltInTok{set} \AttributeTok{{-}gx}\NormalTok{ PATH }\VariableTok{$PATH}\NormalTok{ \textasciitilde{}/android/platform{-}tools}
\end{Highlighting}
\end{Shaded}

\subsection{Problema 2: ``no permissions (missing udev
rules?)''}\label{problema-2-no-permissions-missing-udev-rules}

\textbf{Causa:} Reglas udev faltantes o usuario no en grupo adbusers.

\textbf{Soluciones:}

\begin{Shaded}
\begin{Highlighting}[]
\CommentTok{\# 1. Verificar grupo}
\FunctionTok{groups} \KeywordTok{|} \FunctionTok{grep}\NormalTok{ adbusers}

\CommentTok{\# 2. Agregar a grupo}
\FunctionTok{sudo}\NormalTok{ usermod }\AttributeTok{{-}aG}\NormalTok{ adbusers }\VariableTok{$USER}
\ExtensionTok{newgrp}\NormalTok{ adbusers}

\CommentTok{\# 3. Crear/verificar reglas udev}
\FunctionTok{sudo}\NormalTok{ nano /etc/udev/rules.d/51{-}android.rules}
\CommentTok{\# Ver sección "Configuración Inicial"}

\CommentTok{\# 4. Recargar reglas}
\FunctionTok{sudo}\NormalTok{ udevadm control }\AttributeTok{{-}{-}reload{-}rules}
\FunctionTok{sudo}\NormalTok{ udevadm trigger}

\CommentTok{\# 5. Reconectar dispositivo}
\end{Highlighting}
\end{Shaded}

\subsection{Problema 3: ``device offline'' o
``unauthorized''}\label{problema-3-device-offline-o-unauthorized}

\textbf{Causa:} Dispositivo no autorizado o conexión interrumpida.

\textbf{Soluciones:}

\begin{Shaded}
\begin{Highlighting}[]
\CommentTok{\# 1. Revocar autorizaciones en dispositivo}
\CommentTok{\# Android: Opciones desarrollador → Revocar autorizaciones USB}

\CommentTok{\# 2. Reconectar dispositivo}

\CommentTok{\# 3. Aceptar diálogo de autorización}

\CommentTok{\# 4. Reiniciar servidor adb}
\ExtensionTok{adb}\NormalTok{ kill{-}server}
\ExtensionTok{adb}\NormalTok{ start{-}server}
\ExtensionTok{adb}\NormalTok{ devices}
\end{Highlighting}
\end{Shaded}

\subsection{Problema 4: ``more than one
device/emulator''}\label{problema-4-more-than-one-deviceemulator}

\textbf{Causa:} Múltiples dispositivos conectados sin especificar cuál
usar.

\textbf{Soluciones:}

\begin{Shaded}
\begin{Highlighting}[]
\CommentTok{\# 1. Listar dispositivos}
\ExtensionTok{adb}\NormalTok{ devices}

\CommentTok{\# 2. Especificar por serial}
\ExtensionTok{adb} \AttributeTok{{-}s}\NormalTok{ ABC123456789 shell}

\CommentTok{\# 3. O establecer variable de entorno}
\BuiltInTok{set} \AttributeTok{{-}gx}\NormalTok{ ANDROID\_SERIAL ABC123456789}
\ExtensionTok{adb}\NormalTok{ shell}

\CommentTok{\# 4. Usar {-}d (solo USB) o {-}e (solo emulador)}
\ExtensionTok{adb} \AttributeTok{{-}d}\NormalTok{ shell}
\ExtensionTok{adb} \AttributeTok{{-}e}\NormalTok{ shell}
\end{Highlighting}
\end{Shaded}

\subsection{Problema 5: Conexión WiFi se
desconecta}\label{problema-5-conexiuxf3n-wifi-se-desconecta}

\textbf{Causa:} Red inestable o dispositivo cambia de red.

\textbf{Soluciones:}

\begin{Shaded}
\begin{Highlighting}[]
\CommentTok{\# 1. Verificar que ambos estén en misma red}
\CommentTok{\# Dispositivo: Ajustes → WiFi}
\CommentTok{\# Host: ip addr show}

\CommentTok{\# 2. Reconectar}
\ExtensionTok{adb}\NormalTok{ connect 192.168.1.100:5555}

\CommentTok{\# 3. Si falla, reiniciar desde USB}
\ExtensionTok{adb}\NormalTok{ usb}
\ExtensionTok{adb}\NormalTok{ tcpip 5555}
\ExtensionTok{adb}\NormalTok{ connect 192.168.1.100:5555}

\CommentTok{\# 4. Ping para verificar conectividad}
\FunctionTok{ping}\NormalTok{ 192.168.1.100}
\end{Highlighting}
\end{Shaded}

\subsection{Problema 6: ``daemon not running; starting
now''}\label{problema-6-daemon-not-running-starting-now}

\textbf{Causa:} Servidor ADB no estaba corriendo (normal).

\textbf{No es un error.} ADB inicia automáticamente el servidor.

\begin{Shaded}
\begin{Highlighting}[]
\CommentTok{\# Si quieres iniciar manualmente:}
\ExtensionTok{adb}\NormalTok{ start{-}server}
\end{Highlighting}
\end{Shaded}

\subsection{Problema 7: ``insufficient permissions for
device''}\label{problema-7-insufficient-permissions-for-device}

\textbf{Causa:} Permisos USB incorrectos.

\textbf{Soluciones:}

\begin{Shaded}
\begin{Highlighting}[]
\CommentTok{\# 1. Verificar reglas udev (ver Problema 2)}

\CommentTok{\# 2. Verificar vendor ID del dispositivo}
\ExtensionTok{lsusb}
\CommentTok{\# Buscar tu dispositivo Android}

\CommentTok{\# 3. Agregar vendor ID a reglas udev}
\FunctionTok{sudo}\NormalTok{ nano /etc/udev/rules.d/51{-}android.rules}
\CommentTok{\# SUBSYSTEM=="usb", ATTR\{idVendor\}=="XXXX", MODE="0666", GROUP="adbusers"}

\CommentTok{\# 4. Recargar reglas}
\FunctionTok{sudo}\NormalTok{ udevadm control }\AttributeTok{{-}{-}reload{-}rules}
\end{Highlighting}
\end{Shaded}

\subsection{Problema 8: ``error: protocol fault (no
status)''}\label{problema-8-error-protocol-fault-no-status}

\textbf{Causa:} Comunicación corrupta con dispositivo.

\textbf{Soluciones:}

\begin{Shaded}
\begin{Highlighting}[]
\CommentTok{\# 1. Desconectar y reconectar dispositivo}

\CommentTok{\# 2. Reiniciar servidor ADB}
\ExtensionTok{adb}\NormalTok{ kill{-}server}
\ExtensionTok{adb}\NormalTok{ start{-}server}

\CommentTok{\# 3. Reiniciar dispositivo}
\ExtensionTok{adb}\NormalTok{ reboot}

\CommentTok{\# 4. Como último recurso, reiniciar computadora}
\end{Highlighting}
\end{Shaded}

\section{Automatización con
Scripts}\label{automatizaciuxf3n-con-scripts}

\subsection{Script Fish: Instalador de APK
Mejorado}\label{script-fish-instalador-de-apk-mejorado}

\textbf{Crear
\texttt{\textasciitilde{}/.config/fish/functions/adb-install-smart.fish}:}

\begin{Shaded}
\begin{Highlighting}[]
\NormalTok{function adb{-}install{-}smart {-}{-}description "Instalar APK con opciones inteligentes"}
\NormalTok{    if test (count $argv) {-}lt 1}
\NormalTok{        echo "Uso: adb{-}install{-}smart \textless{}archivo.apk\textgreater{} [opciones]"}
\NormalTok{        echo ""}
\NormalTok{        echo "Opciones:"}
\NormalTok{        echo "  {-}r : Reinstalar manteniendo datos"}
\NormalTok{        echo "  {-}g : Otorgar todos los permisos"}
\NormalTok{        echo "  {-}t : APK de prueba"}
\NormalTok{        echo "  {-}d : Permitir downgrade"}
\NormalTok{        return 1}
\NormalTok{    end}
    
\NormalTok{    set apk $argv[1]}
\NormalTok{    set options}
    
\NormalTok{    \# Parsear opciones}
\NormalTok{    for arg in $argv[2..{-}1]}
\NormalTok{        set options $options $arg}
\NormalTok{    end}
    
\NormalTok{    \# Verificar que el archivo existe}
\NormalTok{    if not test {-}f "$apk"}
\NormalTok{        echo "Error: Archivo no encontrado: $apk"}
\NormalTok{        return 1}
\NormalTok{    end}
    
\NormalTok{    \# Verificar extensión}
\NormalTok{    if not string match {-}q "*.apk" "$apk"}
\NormalTok{        echo "Error: El archivo debe ser .apk"}
\NormalTok{        return 1}
\NormalTok{    end}
    
\NormalTok{    \# Obtener package name}
\NormalTok{    set package (aapt dump badging "$apk" 2\textgreater{}/dev/null | grep package: | awk \textquotesingle{}\{print $2\}\textquotesingle{} | cut {-}d"\textquotesingle{}" {-}f2)}
    
\NormalTok{    echo "Instalando: $apk"}
\NormalTok{    if test {-}n "$package"}
\NormalTok{        echo "Paquete: $package"}
\NormalTok{    end}
    
\NormalTok{    \# Instalar}
\NormalTok{    adb install $options "$apk"}
    
\NormalTok{    if test $status {-}eq 0}
\NormalTok{        echo "✓ Instalación exitosa"}
        
\NormalTok{        \# Si se otorgaron permisos, mostrarlos}
\NormalTok{        if string match {-}q "*{-}g*" "$options"}
\NormalTok{            echo "Permisos otorgados automáticamente"}
\NormalTok{        end}
\NormalTok{    else}
\NormalTok{        echo "✗ Error en la instalación"}
\NormalTok{        return 1}
\NormalTok{    end}
\NormalTok{end}
\end{Highlighting}
\end{Shaded}

\textbf{Usar:}

\begin{Shaded}
\begin{Highlighting}[]
\NormalTok{adb{-}install{-}smart app.apk}
\NormalTok{adb{-}install{-}smart app.apk {-}r {-}g}
\end{Highlighting}
\end{Shaded}

\subsection{Script Bash: Backup
Completo}\label{script-bash-backup-completo}

\textbf{Crear \texttt{\textasciitilde{}/bin/adb-backup-full}:}

\begin{Shaded}
\begin{Highlighting}[]
\CommentTok{\#!/bin/bash}
\CommentTok{\# adb{-}backup{-}full {-} Backup completo del dispositivo}

\BuiltInTok{set} \AttributeTok{{-}e}

\VariableTok{BACKUP\_DIR}\OperatorTok{=}\StringTok{"}\VariableTok{$HOME}\StringTok{/android{-}backups/}\VariableTok{$(}\FunctionTok{date}\NormalTok{ +\%Y\%m\%d\_\%H\%M\%S}\VariableTok{)}\StringTok{"}
\VariableTok{DEVICE\_SERIAL}\OperatorTok{=}\StringTok{""}

\CommentTok{\# Parsear argumentos}
\ControlFlowTok{while} \BuiltInTok{getopts} \StringTok{"s:d:"} \VariableTok{opt}\KeywordTok{;} \ControlFlowTok{do}
    \ControlFlowTok{case} \VariableTok{$opt} \KeywordTok{in}
        \SpecialStringTok{s}\KeywordTok{)} \VariableTok{DEVICE\_SERIAL}\OperatorTok{=}\StringTok{"{-}s }\VariableTok{$OPTARG}\StringTok{"} \ControlFlowTok{;;}
        \SpecialStringTok{d}\KeywordTok{)} \VariableTok{BACKUP\_DIR}\OperatorTok{=}\StringTok{"}\VariableTok{$OPTARG}\StringTok{"} \ControlFlowTok{;;}
        \PreprocessorTok{*}\KeywordTok{)} \BuiltInTok{echo} \StringTok{"Uso: }\VariableTok{$0}\StringTok{ [{-}s serial] [{-}d directorio]"}\KeywordTok{;} \BuiltInTok{exit}\NormalTok{ 1 }\ControlFlowTok{;;}
    \ControlFlowTok{esac}
\ControlFlowTok{done}

\FunctionTok{mkdir} \AttributeTok{{-}p} \StringTok{"}\VariableTok{$BACKUP\_DIR}\StringTok{"}

\BuiltInTok{echo} \StringTok{"==================================="}
\BuiltInTok{echo} \StringTok{"Backup Completo de Dispositivo Android"}
\BuiltInTok{echo} \StringTok{"==================================="}
\BuiltInTok{echo} \StringTok{"Directorio: }\VariableTok{$BACKUP\_DIR}\StringTok{"}
\BuiltInTok{echo} \StringTok{""}

\CommentTok{\# Información del dispositivo}
\BuiltInTok{echo} \StringTok{"1. Capturando información del dispositivo..."}
\ExtensionTok{adb} \VariableTok{$DEVICE\_SERIAL}\NormalTok{ shell getprop }\OperatorTok{\textgreater{}} \StringTok{"}\VariableTok{$BACKUP\_DIR}\StringTok{/device\_info.txt"}
\ExtensionTok{adb} \VariableTok{$DEVICE\_SERIAL}\NormalTok{ shell }\StringTok{"dumpsys battery"} \OperatorTok{\textgreater{}} \StringTok{"}\VariableTok{$BACKUP\_DIR}\StringTok{/battery\_info.txt"}

\CommentTok{\# Lista de aplicaciones}
\BuiltInTok{echo} \StringTok{"2. Listando aplicaciones instaladas..."}
\ExtensionTok{adb} \VariableTok{$DEVICE\_SERIAL}\NormalTok{ shell pm list packages }\AttributeTok{{-}f} \OperatorTok{\textgreater{}} \StringTok{"}\VariableTok{$BACKUP\_DIR}\StringTok{/packages.txt"}
\ExtensionTok{adb} \VariableTok{$DEVICE\_SERIAL}\NormalTok{ shell pm list packages }\AttributeTok{{-}3} \OperatorTok{\textgreater{}} \StringTok{"}\VariableTok{$BACKUP\_DIR}\StringTok{/packages\_user.txt"}

\CommentTok{\# Backup de APKs de usuario}
\BuiltInTok{echo} \StringTok{"3. Respaldando APKs de usuario..."}
\FunctionTok{mkdir} \AttributeTok{{-}p} \StringTok{"}\VariableTok{$BACKUP\_DIR}\StringTok{/apks"}
\VariableTok{USER\_APPS}\OperatorTok{=}\VariableTok{$(}\ExtensionTok{adb} \VariableTok{$DEVICE\_SERIAL}\NormalTok{ shell pm list packages }\AttributeTok{{-}3} \KeywordTok{|} \FunctionTok{cut} \AttributeTok{{-}d:} \AttributeTok{{-}f2}\VariableTok{)}
\ControlFlowTok{for}\NormalTok{ app }\KeywordTok{in} \VariableTok{$USER\_APPS}\KeywordTok{;} \ControlFlowTok{do}
    \BuiltInTok{echo} \StringTok{"  Respaldando: }\VariableTok{$app}\StringTok{"}
    \VariableTok{APK\_PATH}\OperatorTok{=}\VariableTok{$(}\ExtensionTok{adb} \VariableTok{$DEVICE\_SERIAL}\NormalTok{ shell pm path }\StringTok{"}\VariableTok{$app}\StringTok{"} \KeywordTok{|} \FunctionTok{cut} \AttributeTok{{-}d:} \AttributeTok{{-}f2} \KeywordTok{|} \FunctionTok{tr} \AttributeTok{{-}d} \StringTok{\textquotesingle{}\textbackslash{}r\textquotesingle{}}\VariableTok{)}
    \ExtensionTok{adb} \VariableTok{$DEVICE\_SERIAL}\NormalTok{ pull }\StringTok{"}\VariableTok{$APK\_PATH}\StringTok{"} \StringTok{"}\VariableTok{$BACKUP\_DIR}\StringTok{/apks/}\VariableTok{$\{app\}}\StringTok{.apk"} \DecValTok{2}\OperatorTok{\textgreater{}}\NormalTok{/dev/null }\KeywordTok{||} \BuiltInTok{echo} \StringTok{"    Error respaldando }\VariableTok{$app}\StringTok{"}
\ControlFlowTok{done}

\CommentTok{\# Backup de datos}
\BuiltInTok{echo} \StringTok{"4. Respaldando datos de usuario..."}
\ExtensionTok{adb} \VariableTok{$DEVICE\_SERIAL}\NormalTok{ backup }\AttributeTok{{-}f} \StringTok{"}\VariableTok{$BACKUP\_DIR}\StringTok{/data.ab"} \AttributeTok{{-}all} \AttributeTok{{-}apk} \AttributeTok{{-}shared}

\CommentTok{\# Contactos (si accesible)}
\BuiltInTok{echo} \StringTok{"5. Respaldando contactos..."}
\ExtensionTok{adb} \VariableTok{$DEVICE\_SERIAL}\NormalTok{ pull /data/data/com.android.providers.contacts/databases }\StringTok{"}\VariableTok{$BACKUP\_DIR}\StringTok{/contacts/"} \DecValTok{2}\OperatorTok{\textgreater{}}\NormalTok{/dev/null }\KeywordTok{||} \BuiltInTok{echo} \StringTok{"  No accesible (requiere root)"}

\CommentTok{\# Fotos y videos}
\BuiltInTok{echo} \StringTok{"6. Respaldando multimedia..."}
\ExtensionTok{adb} \VariableTok{$DEVICE\_SERIAL}\NormalTok{ pull /sdcard/DCIM }\StringTok{"}\VariableTok{$BACKUP\_DIR}\StringTok{/DCIM/"} \DecValTok{2}\OperatorTok{\textgreater{}}\NormalTok{/dev/null }\KeywordTok{||} \FunctionTok{true}
\ExtensionTok{adb} \VariableTok{$DEVICE\_SERIAL}\NormalTok{ pull /sdcard/Pictures }\StringTok{"}\VariableTok{$BACKUP\_DIR}\StringTok{/Pictures/"} \DecValTok{2}\OperatorTok{\textgreater{}}\NormalTok{/dev/null }\KeywordTok{||} \FunctionTok{true}
\ExtensionTok{adb} \VariableTok{$DEVICE\_SERIAL}\NormalTok{ pull /sdcard/Download }\StringTok{"}\VariableTok{$BACKUP\_DIR}\StringTok{/Download/"} \DecValTok{2}\OperatorTok{\textgreater{}}\NormalTok{/dev/null }\KeywordTok{||} \FunctionTok{true}

\CommentTok{\# Crear archivo de metadatos}
\BuiltInTok{echo} \StringTok{"7. Creando metadatos..."}
\FunctionTok{cat} \OperatorTok{\textgreater{}} \StringTok{"}\VariableTok{$BACKUP\_DIR}\StringTok{/README.txt"} \OperatorTok{\textless{}\textless{} EOF}
\StringTok{Backup de Dispositivo Android}
\StringTok{Fecha: }\VariableTok{$(}\FunctionTok{date}\VariableTok{)}
\StringTok{Dispositivo: }\VariableTok{$(}\ExtensionTok{adb} \VariableTok{$DEVICE\_SERIAL}\NormalTok{ shell getprop ro.product.model}\VariableTok{)}
\StringTok{Android: }\VariableTok{$(}\ExtensionTok{adb} \VariableTok{$DEVICE\_SERIAL}\NormalTok{ shell getprop ro.build.version.release}\VariableTok{)}
\StringTok{Serial: }\VariableTok{$(}\ExtensionTok{adb} \VariableTok{$DEVICE\_SERIAL}\NormalTok{ get{-}serialno}\VariableTok{)}

\StringTok{Contenido:}
\StringTok{{-} device\_info.txt: Propiedades del sistema}
\StringTok{{-} packages.txt: Lista de todos los paquetes}
\StringTok{{-} packages\_user.txt: Solo paquetes de usuario}
\StringTok{{-} apks/: APKs de aplicaciones de usuario}
\StringTok{{-} data.ab: Backup de datos (usar \textquotesingle{}adb restore data.ab\textquotesingle{})}
\StringTok{{-} DCIM/, Pictures/, Download/: Multimedia}
\OperatorTok{EOF}

\BuiltInTok{echo} \StringTok{""}
\BuiltInTok{echo} \StringTok{"==================================="}
\BuiltInTok{echo} \StringTok{"Backup completado exitosamente"}
\BuiltInTok{echo} \StringTok{"Ubicación: }\VariableTok{$BACKUP\_DIR}\StringTok{"}
\BuiltInTok{echo} \StringTok{"==================================="}
\end{Highlighting}
\end{Shaded}

\textbf{Hacer ejecutable y usar:}

\begin{Shaded}
\begin{Highlighting}[]
\FunctionTok{chmod}\NormalTok{ +x \textasciitilde{}/bin/adb{-}backup{-}full}
\ExtensionTok{adb{-}backup{-}full}
\ExtensionTok{adb{-}backup{-}full} \AttributeTok{{-}s}\NormalTok{ ABC123456789}
\ExtensionTok{adb{-}backup{-}full} \AttributeTok{{-}d}\NormalTok{ \textasciitilde{}/mis{-}backups/backup{-}especial}
\end{Highlighting}
\end{Shaded}

\subsection{Alias Útiles para Fish}\label{alias-uxfatiles-para-fish}

\textbf{Agregar a \texttt{\textasciitilde{}/.config/fish/config.fish}:}

\begin{Shaded}
\begin{Highlighting}[]
\NormalTok{\# =============================================================================}
\NormalTok{\# ALIAS DE ADB}
\NormalTok{\# =============================================================================}

\NormalTok{\# Básicos}
\NormalTok{alias adb{-}devices="adb devices {-}l"}
\NormalTok{alias adb{-}restart="adb kill{-}server \&\& adb start{-}server"}
\NormalTok{alias adb{-}logcat="adb logcat {-}v time"}

\NormalTok{\# Dispositivo}
\NormalTok{alias adb{-}model="adb shell getprop ro.product.model"}
\NormalTok{alias adb{-}android="adb shell getprop ro.build.version.release"}
\NormalTok{alias adb{-}battery="adb shell dumpsys battery"}
\NormalTok{alias adb{-}wifi="adb shell svc wifi enable"}
\NormalTok{alias adb{-}wifi{-}off="adb shell svc wifi disable"}

\NormalTok{\# Aplicaciones}
\NormalTok{alias adb{-}apps="adb shell pm list packages {-}3"}
\NormalTok{alias adb{-}apps{-}all="adb shell pm list packages"}
\NormalTok{alias adb{-}uninstall="adb uninstall"}
\NormalTok{alias adb{-}clear="adb shell pm clear"}

\NormalTok{\# Archivos}
\NormalTok{alias adb{-}pull{-}dcim="adb pull /sdcard/DCIM/ ./"}
\NormalTok{alias adb{-}pull{-}download="adb pull /sdcard/Download/ ./"}
\NormalTok{alias adb{-}ls="adb shell ls {-}la"}

\NormalTok{\# Captura}
\NormalTok{alias adb{-}screenshot="adb exec{-}out screencap {-}p \textgreater{} screenshot\_(date +\%Y\%m\%d\_\%H\%M\%S).png"}
\NormalTok{alias adb{-}record="adb shell screenrecord /sdcard/video.mp4"}

\NormalTok{\# Desarrollo}
\NormalTok{alias adb{-}install{-}debug="adb install {-}r {-}t {-}g"}
\NormalTok{alias adb{-}logs="adb logcat {-}v time | grep {-}i"}

\NormalTok{\# Conexión}
\NormalTok{alias adb{-}wifi{-}connect="adb tcpip 5555"}
\NormalTok{alias adb{-}usb="adb usb"}

\NormalTok{\# Reinicio}
\NormalTok{alias adb{-}reboot="adb reboot"}
\NormalTok{alias adb{-}recovery="adb reboot recovery"}
\NormalTok{alias adb{-}bootloader="adb reboot bootloader"}
\end{Highlighting}
\end{Shaded}

\textbf{Recargar:}

\begin{Shaded}
\begin{Highlighting}[]
\NormalTok{source \textasciitilde{}/.config/fish/config.fish}
\end{Highlighting}
\end{Shaded}

\section{Referencia Rápida}\label{referencia-ruxe1pida}

\subsection{Comandos Esenciales}\label{comandos-esenciales}

\begin{Shaded}
\begin{Highlighting}[]
\CommentTok{\# SERVIDOR}
\ExtensionTok{adb}\NormalTok{ start{-}server              }\CommentTok{\# Iniciar servidor}
\ExtensionTok{adb}\NormalTok{ kill{-}server               }\CommentTok{\# Detener servidor}
\ExtensionTok{adb}\NormalTok{ version                   }\CommentTok{\# Versión}

\CommentTok{\# DISPOSITIVOS}
\ExtensionTok{adb}\NormalTok{ devices                   }\CommentTok{\# Listar dispositivos}
\ExtensionTok{adb}\NormalTok{ devices }\AttributeTok{{-}l}                \CommentTok{\# Con detalles}
\ExtensionTok{adb} \AttributeTok{{-}s}\NormalTok{ SERIAL comando         }\CommentTok{\# Especificar dispositivo}
\ExtensionTok{adb} \AttributeTok{{-}d}\NormalTok{ comando                }\CommentTok{\# Solo USB}
\ExtensionTok{adb} \AttributeTok{{-}e}\NormalTok{ comando                }\CommentTok{\# Solo emulador}

\CommentTok{\# CONEXIÓN}
\ExtensionTok{adb}\NormalTok{ tcpip 5555                }\CommentTok{\# Habilitar WiFi}
\ExtensionTok{adb}\NormalTok{ connect IP:5555           }\CommentTok{\# Conectar WiFi}
\ExtensionTok{adb}\NormalTok{ disconnect                }\CommentTok{\# Desconectar WiFi}
\ExtensionTok{adb}\NormalTok{ usb                       }\CommentTok{\# Volver a USB}

\CommentTok{\# INFORMACIÓN}
\ExtensionTok{adb}\NormalTok{ shell getprop             }\CommentTok{\# Propiedades}
\ExtensionTok{adb}\NormalTok{ get{-}serialno              }\CommentTok{\# Número de serie}
\ExtensionTok{adb}\NormalTok{ get{-}state                 }\CommentTok{\# Estado}

\CommentTok{\# APLICACIONES}
\ExtensionTok{adb}\NormalTok{ install app.apk           }\CommentTok{\# Instalar}
\ExtensionTok{adb}\NormalTok{ install }\AttributeTok{{-}r}\NormalTok{ app.apk        }\CommentTok{\# Reinstalar}
\ExtensionTok{adb}\NormalTok{ uninstall PACKAGE         }\CommentTok{\# Desinstalar}
\ExtensionTok{adb}\NormalTok{ shell pm list packages    }\CommentTok{\# Listar apps}
\ExtensionTok{adb}\NormalTok{ shell pm clear PACKAGE    }\CommentTok{\# Limpiar datos}

\CommentTok{\# ARCHIVOS}
\ExtensionTok{adb}\NormalTok{ push LOCAL REMOTE         }\CommentTok{\# Copiar a dispositivo}
\ExtensionTok{adb}\NormalTok{ pull REMOTE LOCAL         }\CommentTok{\# Copiar de dispositivo}
\ExtensionTok{adb}\NormalTok{ shell ls /sdcard          }\CommentTok{\# Listar archivos}

\CommentTok{\# SHELL}
\ExtensionTok{adb}\NormalTok{ shell                     }\CommentTok{\# Shell interactivo}
\ExtensionTok{adb}\NormalTok{ shell COMANDO             }\CommentTok{\# Comando único}

\CommentTok{\# CAPTURA}
\ExtensionTok{adb}\NormalTok{ exec{-}out screencap }\AttributeTok{{-}p} \OperatorTok{\textgreater{}}\NormalTok{ screenshot.png}
\ExtensionTok{adb}\NormalTok{ shell screenrecord /sdcard/video.mp4}

\CommentTok{\# LOGS}
\ExtensionTok{adb}\NormalTok{ logcat                    }\CommentTok{\# Ver logs}
\ExtensionTok{adb}\NormalTok{ logcat }\AttributeTok{{-}c}                 \CommentTok{\# Limpiar logs}
\ExtensionTok{adb}\NormalTok{ bugreport                 }\CommentTok{\# Reporte de bugs}

\CommentTok{\# REINICIO}
\ExtensionTok{adb}\NormalTok{ reboot                    }\CommentTok{\# Reiniciar}
\ExtensionTok{adb}\NormalTok{ reboot recovery           }\CommentTok{\# Recovery}
\ExtensionTok{adb}\NormalTok{ reboot bootloader         }\CommentTok{\# Bootloader}
\end{Highlighting}
\end{Shaded}

\subsection{Atajos de Teclado en
Terminal}\label{atajos-de-teclado-en-terminal}

\begin{longtable}[]{@{}ll@{}}
\toprule\noalign{}
Atajo & Acción \\
\midrule\noalign{}
\endhead
\bottomrule\noalign{}
\endlastfoot
\texttt{Ctrl+C} & Detener comando actual \\
\texttt{Ctrl+D} & Salir de shell adb \\
\texttt{Ctrl+Z} & Suspender comando \\
\texttt{Ctrl+L} & Limpiar pantalla \\
\texttt{↑} / \texttt{↓} & Historial de comandos \\
\end{longtable}

\subsection{Variables de Entorno
Útiles}\label{variables-de-entorno-uxfatiles}

\begin{Shaded}
\begin{Highlighting}[]
\NormalTok{\# Fish}
\NormalTok{set {-}gx ANDROID\_SERIAL ABC123456789}
\NormalTok{set {-}gx ANDROID\_HOME /opt/android{-}sdk}
\NormalTok{set {-}gx ADB\_TRACE all  \# Debug verbose}

\NormalTok{\# Bash/Zsh}
\NormalTok{export ANDROID\_SERIAL=ABC123456789}
\NormalTok{export ANDROID\_HOME=/opt/android{-}sdk}
\NormalTok{export ADB\_TRACE=all}
\end{Highlighting}
\end{Shaded}

\subsection{Locations Importantes}\label{locations-importantes}

\begin{verbatim}
/sdcard/              # Almacenamiento interno (equivale a /storage/emulated/0/)
/sdcard/DCIM/         # Fotos de cámara
/sdcard/Download/     # Descargas
/sdcard/Pictures/     # Imágenes
/sdcard/Music/        # Música
/sdcard/Movies/       # Videos
/data/data/           # Datos de apps (requiere root)
/system/app/          # Apps del sistema (requiere root)
/data/app/            # Apps de usuario (requiere root)
\end{verbatim}

\subsection{Intents Comunes}\label{intents-comunes}

\begin{Shaded}
\begin{Highlighting}[]
\CommentTok{\# Configuración}
\ExtensionTok{android.settings.SETTINGS}
\ExtensionTok{android.settings.WIFI\_SETTINGS}
\ExtensionTok{android.settings.BLUETOOTH\_SETTINGS}
\ExtensionTok{android.settings.APPLICATION\_SETTINGS}

\CommentTok{\# Acciones}
\ExtensionTok{android.intent.action.VIEW}
\ExtensionTok{android.intent.action.CALL}
\ExtensionTok{android.intent.action.SENDTO}
\ExtensionTok{android.intent.action.SEND}
\ExtensionTok{android.media.action.IMAGE\_CAPTURE}
\end{Highlighting}
\end{Shaded}

\section{Recursos Adicionales}\label{recursos-adicionales}

\subsection{Documentación Oficial}\label{documentaciuxf3n-oficial}

\begin{itemize}
\tightlist
\item
  \textbf{ADB Official:}
  https://developer.android.com/studio/command-line/adb
\item
  \textbf{Platform Tools:}
  https://developer.android.com/studio/releases/platform-tools
\item
  \textbf{Android Debug:} https://developer.android.com/studio/debug
\end{itemize}

\subsection{Comunidad}\label{comunidad}

\begin{itemize}
\tightlist
\item
  \textbf{XDA Developers:} https://www.xda-developers.com/
\item
  \textbf{Stack Overflow:}
  https://stackoverflow.com/questions/tagged/adb
\item
  \textbf{Android Developers:} https://developer.android.com/
\end{itemize}

\subsection{Herramientas
Complementarias}\label{herramientas-complementarias}

\begin{Shaded}
\begin{Highlighting}[]
\CommentTok{\# En Arch Linux}
\FunctionTok{sudo}\NormalTok{ pacman }\AttributeTok{{-}S}\NormalTok{ android{-}tools android{-}udev scrcpy}

\CommentTok{\# Otras herramientas útiles}
\ExtensionTok{yay} \AttributeTok{{-}S}\NormalTok{ apktool android{-}sdk{-}platform{-}tools}
\end{Highlighting}
\end{Shaded}

\section{Mejores Prácticas}\label{mejores-pruxe1cticas}

\subsection{Seguridad}\label{seguridad}

\begin{enumerate}
\def\labelenumi{\arabic{enumi}.}
\item
  \textbf{Depuración USB solo cuando sea necesario}

  \begin{itemize}
  \tightlist
  \item
    Desactivar cuando no uses ADB
  \item
    No conectar a computadoras desconocidas
  \end{itemize}
\item
  \textbf{Autorizar solo computadoras confiables}

  \begin{itemize}
  \tightlist
  \item
    Verificar huella digital RSA
  \item
    Revocar autorizaciones periódicamente
  \end{itemize}
\item
  \textbf{Conexión WiFi segura}

  \begin{itemize}
  \tightlist
  \item
    Usar solo en redes confiables
  \item
    Desactivar cuando no la uses
  \item
    No usar en redes públicas
  \end{itemize}
\end{enumerate}

\subsection{Rendimiento}\label{rendimiento}

\begin{enumerate}
\def\labelenumi{\arabic{enumi}.}
\item
  \textbf{Cerrar servidor cuando no lo uses}

\begin{Shaded}
\begin{Highlighting}[]
\ExtensionTok{adb}\NormalTok{ kill{-}server}
\end{Highlighting}
\end{Shaded}
\item
  \textbf{Usar comandos específicos}

\begin{Shaded}
\begin{Highlighting}[]
\CommentTok{\# Bien}
\ExtensionTok{adb} \AttributeTok{{-}s}\NormalTok{ ABC123 shell pm list packages }\AttributeTok{{-}3}

\CommentTok{\# Evitar}
\ExtensionTok{adb}\NormalTok{ shell pm list packages }\KeywordTok{|} \FunctionTok{grep}\NormalTok{ algo}
\end{Highlighting}
\end{Shaded}
\item
  \textbf{Transferencias grandes por USB}

  \begin{itemize}
  \tightlist
  \item
    WiFi es más lento para archivos grandes
  \item
    USB 3.0 es significativamente más rápido
  \end{itemize}
\end{enumerate}

\subsection{Desarrollo}\label{desarrollo}

\begin{enumerate}
\def\labelenumi{\arabic{enumi}.}
\item
  \textbf{Automatizar tareas repetitivas}

  \begin{itemize}
  \tightlist
  \item
    Crear scripts Fish/Bash
  \item
    Usar alias para comandos comunes
  \end{itemize}
\item
  \textbf{Mantener logs organizados}

\begin{Shaded}
\begin{Highlighting}[]
\ExtensionTok{adb}\NormalTok{ logcat }\OperatorTok{\textgreater{}}\NormalTok{ logs/app\_}\VariableTok{$(}\FunctionTok{date}\NormalTok{ +\%Y\%m\%d\_\%H\%M\%S}\VariableTok{)}\NormalTok{.log}
\end{Highlighting}
\end{Shaded}
\item
  \textbf{Backup regular}

  \begin{itemize}
  \tightlist
  \item
    Antes de pruebas destructivas
  \item
    Usar scripts automatizados
  \end{itemize}
\end{enumerate}

\section{Publicaciones Similares}\label{publicaciones-similares}

Si te interesó este artículo, te recomendamos que explores otros blogs y
recursos relacionados que pueden ampliar tus conocimientos. Aquí te dejo
algunas sugerencias:

\begin{enumerate}
\def\labelenumi{\arabic{enumi}.}
\tightlist
\item
  \href{https://chaska-x.netlify.app/operating-system/2017-05-21-comandos-de-informacion-windows/index.pdf}{\faIcon{file-pdf}}
  \href{https://chaska-x.netlify.app/operating-system/2017-05-21-comandos-de-informacion-windows}{Comandos
  De Informacion Windows}
\item
  \href{https://chaska-x.netlify.app/operating-system/2019-06-19-adb/index.pdf}{\faIcon{file-pdf}}
  \href{https://chaska-x.netlify.app/operating-system/2019-06-19-adb}{Adb}
\item
  \href{https://chaska-x.netlify.app/operating-system/2021-08-17-limpieza-y-optimizacion-de-pc/index.pdf}{\faIcon{file-pdf}}
  \href{https://chaska-x.netlify.app/operating-system/2021-08-17-limpieza-y-optimizacion-de-pc}{Limpieza
  Y Optimizacion De Pc}
\item
  \href{https://chaska-x.netlify.app/operating-system/2021-10-21-usando-apk-en-windown-11/index.pdf}{\faIcon{file-pdf}}
  \href{https://chaska-x.netlify.app/operating-system/2021-10-21-usando-apk-en-windown-11}{Usando
  Apk En Windown 11}
\item
  \href{https://chaska-x.netlify.app/operating-system/2022-05-12-gestionar-versiones-de-jdk-en-kubuntu/index.pdf}{\faIcon{file-pdf}}
  \href{https://chaska-x.netlify.app/operating-system/2022-05-12-gestionar-versiones-de-jdk-en-kubuntu}{Gestionar
  Versiones De Jdk En Kubuntu}
\item
  \href{https://chaska-x.netlify.app/operating-system/2022-07-21-instalar-tor-browser/index.pdf}{\faIcon{file-pdf}}
  \href{https://chaska-x.netlify.app/operating-system/2022-07-21-instalar-tor-browser}{Instalar
  Tor Browser}
\item
  \href{https://chaska-x.netlify.app/operating-system/2022-08-14-crear-enlaces-duros-o-hard-link-en-linux/index.pdf}{\faIcon{file-pdf}}
  \href{https://chaska-x.netlify.app/operating-system/2022-08-14-crear-enlaces-duros-o-hard-link-en-linux}{Crear
  Enlaces Duros O Hard Link En Linux}
\item
  \href{https://chaska-x.netlify.app/operating-system/2022-09-27-comandos-vim/index.pdf}{\faIcon{file-pdf}}
  \href{https://chaska-x.netlify.app/operating-system/2022-09-27-comandos-vim}{Comandos
  Vim}
\item
  \href{https://chaska-x.netlify.app/operating-system/2023-02-16-guia-de-git-y-github/index.pdf}{\faIcon{file-pdf}}
  \href{https://chaska-x.netlify.app/operating-system/2023-02-16-guia-de-git-y-github}{Guia
  De Git Y Github}
\item
  \href{https://chaska-x.netlify.app/operating-system/2023-05-02-00-primeros-pasos-en-linux/index.pdf}{\faIcon{file-pdf}}
  \href{https://chaska-x.netlify.app/operating-system/2023-05-02-00-primeros-pasos-en-linux}{00
  Primeros Pasos En Linux}
\item
  \href{https://chaska-x.netlify.app/operating-system/2023-06-17-01-introduccion-linux/index.pdf}{\faIcon{file-pdf}}
  \href{https://chaska-x.netlify.app/operating-system/2023-06-17-01-introduccion-linux}{01
  Introduccion Linux}
\item
  \href{https://chaska-x.netlify.app/operating-system/2023-06-18-02-distribuciones-linux/index.pdf}{\faIcon{file-pdf}}
  \href{https://chaska-x.netlify.app/operating-system/2023-06-18-02-distribuciones-linux}{02
  Distribuciones Linux}
\item
  \href{https://chaska-x.netlify.app/operating-system/2023-06-19-03-instalacion-linux/index.pdf}{\faIcon{file-pdf}}
  \href{https://chaska-x.netlify.app/operating-system/2023-06-19-03-instalacion-linux}{03
  Instalacion Linux}
\item
  \href{https://chaska-x.netlify.app/operating-system/2023-06-20-04-administracion-particiones-volumenes/index.pdf}{\faIcon{file-pdf}}
  \href{https://chaska-x.netlify.app/operating-system/2023-06-20-04-administracion-particiones-volumenes}{04
  Administracion Particiones Volumenes}
\item
  \href{https://chaska-x.netlify.app/operating-system/2023-07-01-atajos-de-teclado-y-comandos-para-usar-vim/index.pdf}{\faIcon{file-pdf}}
  \href{https://chaska-x.netlify.app/operating-system/2023-07-01-atajos-de-teclado-y-comandos-para-usar-vim}{Atajos
  De Teclado Y Comandos Para Usar Vim}
\item
  \href{https://chaska-x.netlify.app/operating-system/2024-07-15-instalando-specitify/index.pdf}{\faIcon{file-pdf}}
  \href{https://chaska-x.netlify.app/operating-system/2024-07-15-instalando-specitify}{Instalando
  Specitify}
\item
  \href{https://chaska-x.netlify.app/operating-system/2025-07-10-gestiona-tus-dotfiles-con-gnu-stow/index.pdf}{\faIcon{file-pdf}}
  \href{https://chaska-x.netlify.app/operating-system/2025-07-10-gestiona-tus-dotfiles-con-gnu-stow}{Gestiona
  Tus Dotfiles Con Gnu Stow}
\end{enumerate}

Esperamos que encuentres estas publicaciones igualmente interesantes y
útiles. ¡Disfruta de la lectura!






\end{document}
